\documentclass[11pt]{article}
\usepackage[T1]{fontenc}
\usepackage{lmodern,amsmath,amssymb,amsthm,mathtools}
\usepackage[margin=0.88in]{geometry}
\usepackage{microtype,booktabs,longtable,array,xcolor}
\usepackage{enumitem,fancyhdr,xurl,hyperref}
\hypersetup{colorlinks=true,linkcolor=blue!45!black,urlcolor=blue!45!black,citecolor=blue!45!black,
 pdftitle={Joint arithmetic certificates for ZFC consistency},pdfauthor={AI-assisted research prepared for Vladimir Reshetnikov}}
\pagestyle{fancy}\fancyhf{}\fancyhead[L]{\small Joint arithmetic certificates for ZFC consistency}\fancyhead[R]{\small 25 October 2026}\fancyfoot[C]{\thepage}
\setlength{\headheight}{14pt}
\setlength{\parskip}{4pt}\setlength{\parindent}{0pt}
\setlist{nosep,leftmargin=*}
\allowdisplaybreaks[4]
\newtheorem{theorem}{Theorem}[section]
\newtheorem{lemma}[theorem]{Lemma}
\newtheorem{proposition}[theorem]{Proposition}
\newtheorem{corollary}[theorem]{Corollary}
\theoremstyle{definition}\newtheorem{definition}[theorem]{Definition}
\newcommand{\Z}{\mathbb Z}\newcommand{\N}{\mathbb N}
\newcommand{\Con}{\operatorname{Con}}\newcommand{\sgn}{\operatorname{sgn}}
\newcommand{\code}[1]{\texttt{\detokenize{#1}}}
\newcommand{\pair}[2]{\bigl(X_{#1}(#2),Y_{#1}(#2)\bigr)}

\title{\textbf{Joint arithmetic certificates for ZFC consistency}\\[6pt]
\large Signed Pell indices and exact powers in one polynomial condition}
\author{Prepared for Vladimir Reshetnikov and the ProveIt repository\\
\small AI-assisted mathematical research, with separate internal reviews}
\date{25 October 2026}

\begin{document}
\maketitle
\begin{abstract}
We construct explicit ordinary integer polynomials whose lack of integer zeros
is provably equivalent in ZFC to its syntactic consistency. The new arithmetic
step combines two formerly separate tasks: authenticating the signs and indices
of Pell solutions, and recovering seventeen exact powers at five bases. A
single rational sum is forced to be an integer, then represented by a same-sign
sum of bounded errors of absolute value below one. Its polynomial numerator
replaces eleven private coordinates by one. Zero bases are handled before the
source's sign bootstrap, and the power congruences use bounded classified
indices before the requested indices are authenticated. This avoids a circular
application of the canonical-power lemma.

Applied to the pinned incoming 216-variable quartic, the construction gives a
literal 206-variable polynomial of exact degree 26. A 207-variable degree-24
intermediate and several partitions expose other tradeoffs. Certified monic
substitution then reduces the variable counts substantially; the measured
family is listed in Section~\ref{sec:results}. Every final equation is supplied
as a finite arithmetic circuit with literal integer constants, together with
reproducible construction and independent exact checks. The calibration to
ZFC refutations is the explicitly identified inherited C222 calibration.
No global minimum, literature-wide priority, new proof of consistency, or
newly compiled end-to-end proof-assistant derivation is asserted.
\end{abstract}
\tableofcontents

\section{The outcome and the comparison being made}
\label{sec:scope}

The requested target is a \emph{fixed} polynomial $P\in\Z[x_1,\ldots,x_n]$
for which
\begin{equation}\label{eq:target}
\mathrm{ZFC}\vdash
 \left(\neg\exists\boldsymbol{x}\in\Z^n\;P(\boldsymbol{x})=0\right)
 \longleftrightarrow\Con(\mathrm{ZFC}).
\end{equation}
Here $\Con(\mathrm{ZFC})$ means that the effective syntactic presentation
named by the inherited calibration has no finite proof of contradiction.
The equivalence is proved without assuming either consistency or
inconsistency. Its two constructive directions translate a zero into a
refutation and a refutation into a zero.

There is no single established meaning of ``simplest'' across variable count,
degree, coefficient size, expanded support, and evaluation cost. The repository
already contains equations in very few variables using large fixed coefficient
encodings. In particular its 37-integer quartic belongs to a different
coefficient architecture from the direct construction considered here.
The direct branch instead encodes the numerical calculus C222 with small
literal constants; it does not specialize Jones's universal polynomial or
encode a run of K215. We improve and explicitly compare that branch.

The comparison snapshot is ProveIt commit
\code{315154b670d7bdb1679554734e417410c9b8187b}.
The exact parent is the $P_{216}$ equation in the incoming archive
\code{ConZFC_Pooled_Projection_2026-10-25.zip}~\cite{parent}.
Its SHA-256 is
\begin{quote}\small\ttfamily\sloppy
a91308bd0c82182ea15080763cb23cb4fa0aa45a812eabb7a6cbcd21c0b95a06.
\end{quote}
It has 216 unrestricted integer variables, 153 quadratic residuals,
15,158 collected quartic monomials, and coefficient height 334,952.
The archive's finite verifier was rerun successfully for this research.
The retained original archive contains its ordinary arithmetic proof,
construction, algebra certificates, and named calibration sources.

Three earlier ingredients must receive precise credit. Source 158 developed
same-base power pooling. Source 168 combined two consecutive bases and
suggested using base weights to align signs in larger pools~\cite{pool168}.
Source 152 developed the monic substitution method used in our subsequent
variable elimination~\cite{elim152}. The new contribution here is the
complete simultaneous treatment of the signed-index pool and all power pools,
including zero-base tolerance, its acyclic application to the actual source,
and the emitted equations and certified reductions. We have not established
that this general idea is absent from all mathematical literature.

For context, universal constructions with very few integer unknowns have
also been studied independently, including a recent formalized analysis of
eleven-unknown universal pairs~\cite{bayer}. Their variable-degree bounds
address a different optimization problem and do not imply minimality of any
particular consistency equation.

\section{Pell coordinates, signed indices, and the source frame}
\label{sec:pell}

For an integer $A$, define the canonical Pell polynomials by
\[
 X_0(A)=1,\qquad Y_0(A)=0,
\]
\[
 X_{k+1}=AX_k+(A^2-1)Y_k,\qquad Y_{k+1}=X_k+AY_k.
\]
For negative indices set $X_{-k}=X_k$ and $Y_{-k}=-Y_k$.
Then $X_k^2-(A^2-1)Y_k^2=1$. This notation permits negative $A$ and
negative indices; neither raw coordinate is assumed positive.

\begin{lemma}[Two simultaneous residue identities]\label{lem:residues}
Let $|A|>1$, $L=A-1$, and $T_n=2An-n^2-1$. Then for every signed integer
index $k$,
\[
 X_k(A)\equiv1\pmod L,\qquad Y_k(A)\equiv k\pmod L.
\]
If $n\ne0$, then $\gcd(n,T_n)=1$ and
\begin{equation}\label{eq:pellpower}
 X_k(A)+(n-A)Y_k(A)\equiv n^k\pmod{T_n},
\end{equation}
where a negative power denotes the inverse of the unit $n$ in the residue
ring.
\end{lemma}
\begin{proof}
The first two congruences follow by induction from the recurrences modulo
$A-1$, and by the definitions at negative indices. Also $T_n\equiv-1\pmod n$.
For the power congruence, work in the ring
$\Z[u]/(u^2-(A^2-1))$. The element $A+u$ is a unit with inverse $A-u$,
and its signed $k$th power is $X_k+uY_k$. The substitution $u\mapsto n-A$
defines a map to $\Z/T_n\Z$, since
$(n-A)^2-(A^2-1)=-T_n$. It sends $A+u$ to $n$, proving the claim
for positive, zero, and negative $k$ at once.
\end{proof}

The source has fourteen principal norm-one pairs and seventeen power
requests. Some requests share a pair. Its actual parameter is
$A=x_{134}+1$, and its explicit square frame is
\begin{equation}\label{eq:frame}
 W=2+\sum_{a\in\mathcal A}a^2,
\end{equation}
where the 36 affine atoms in $\mathcal A$ include every requested base,
index, and output, together with the nominal anchor index $h$.
The horizon is $t=16+4d^2$. Consequently
\begin{equation}\label{eq:bounds}
 W\ge258,\qquad
 b:=\lfloor\sqrt{W-1}\rfloor\ge16,\qquad
 |n_j|,|m_r|,|k_i|,|h|\le b.
\end{equation}
The symbol $b$ occurs only in the proof; it is not an extra variable or
a floor operation in the polynomial.

The final parent row basis combines the frame with other equations. It
would be incorrect to interpret its row 129 alone as \eqref{eq:frame}.
The independent port audit reconstructs exactly
\begin{equation}\label{eq:framecomb}
 W-2-\sum_{a\in\mathcal A}a^2
 =22Q_4+Q_7-5Q_{20}+Q_{21}-Q_{46}+Q_{129}+Q_{146}.
\end{equation}
Every row on the right is retained in both new constructions. All row and
coordinate numbers in this article refer to the original 216-coordinate
parent, unless explicitly stated otherwise.

\subsection{The large parameter is available before power recovery}

Rows 92--94, also retained, give
\[
 U=W(2z_0-1),\qquad V=(W+1)U,\qquad A^2+U^2-V^2=1.
\]
Thus $A^2-W(W+2)U^2=1$, and $U$ is nonzero and divisible by $W$.
The elementary classification of solutions to this Pell equation gives
$|A|=X_m(W+1)$ and $|U|=Y_m(W+1)$ for some $m\ge1$.
The recurrence modulo $W$ gives $Y_m(W+1)\equiv m\pmod W$, so $W\mid m$.
The standard growth estimate $X_m(a)\ge a^m$ for $a\ge2$ now yields
\begin{equation}\label{eq:largeA}
 |A|\ge(W+1)^W>W^W.
\end{equation}
This argument uses neither an exact-power certificate nor the orientation
quotient that will be deleted.

\subsection{Bounded classified indices precede orientation}
\label{sec:rawindex}

The second needed ingredient is stronger than a bound on the requested
indices: it bounds the actual indices of the raw Pell solutions before
those indices have been authenticated. Write $a=|A|$. The retained squared
anchor and comparison equations give an anchor ordinate $C$ and
\begin{equation}\label{eq:strictcap}
 C^2=1+h^2+\sum_{i\ne\mathrm{anchor}}Y_i^2
          +u^2+v^2+w(w+1).
\end{equation}
All coordinates are integers, so the last three terms are nonnegative.
In particular $C\ne0$ and $|C|>|h|$.

The signed anchor theorem established in the inherited proof~\cite{anchor135}
applies before the global orientation row. Its Pell classification and
comparison residues imply, for the anchor's positive classified index $j$,
\[
 1\le j\le |C|,\qquad j\equiv\pm h\pmod{2|C|}.
\]
The case $h=0$ is impossible. Since $|h|<|C|$, the congruence gives
$j=|h|$, hence $|C|=Y_{|h|}(a)$. Equation~\eqref{eq:strictcap} and
monotonicity give smaller classified indices for all other pairs.

Therefore each raw pair has a representation
\begin{equation}\label{eq:rawclass}
 (X_i,Y_i)=\varepsilon_i\pair{\nu_i}{A},\qquad
 \varepsilon_i\in\{1,-1\},\quad \nu_i\in\Z,\quad |\nu_i|\le b.
\end{equation}
This representation includes $(1,0)$ and $(-1,0)$. The signed parameter
and independent original signs are absorbed into $\varepsilon_i$ and
the sign of $\nu_i$.

The raw-index bound is an inherited unbounded theorem, not inferred from
testing examples. Its application here was independently checked against
the literal current anchor, norm, cap, and comparison rows. The audit
records their exact integer row combinations and verifies that none uses
the deleted orientation row 142 or any deleted power row. This dependency
order is the reason the next joint argument is legitimate.

\section{A single certificate for all power requests}
\label{sec:powerpool}

First suppose the raw pairs already have their requested signed indices.
For at most five base groups $I_j$, with at most seventeen requests in
total, put
\begin{align}
 T_j&=2An_j-n_j^2-1,\label{eq:T}\\
 H_j&=\sum_{r\in I_j}Y_{i(r)}(m_r-X_{i(r)}),\label{eq:H}\\
 S_j&=\sum_{r\in I_j}\left((m_r-X_{i(r)}+Z_{i(r)})^2-Y_{i(r)}^2\right),\label{eq:S}\\
 E_j&=S_j-2n_jH_j,\qquad Z_i=AY_i.\label{eq:E}
\end{align}
The symbols in these definitions are polynomial abbreviations, not extra
quantified variables. In the literal source the $Z_i$ values are already
represented by retained coordinates and their defining equations.

\begin{theorem}[Oriented pooling with zero-base tolerance]\label{thm:power}
Assume \eqref{eq:bounds}, \eqref{eq:largeA}, canonical signed pairs at
the requested indices, and $Z_i=AY_i$. Then there is an integer $z$ with
\begin{equation}\label{eq:R}
 \mathcal R=
 \sum_{j=0}^4 n_jE_j\prod_{\ell\ne j}T_\ell
       -z\prod_{j=0}^4T_j=0
\end{equation}
if and only if every request at every nonzero base has the exact value
$m_r=n_j^{k_{i(r)}}$. At a base with $|n_j|\ge2$, the requested index must be
nonnegative. Unit bases permit signed indices. Zero bases impose no power
condition.
\end{theorem}

\begin{proof}
For a zero base, $T_j=-1$. For a nonzero base,
\begin{equation}\label{eq:Tbounds}
 |T_j|\ge2|A|-W>W^W,
 \qquad\sgn(T_j)=\sgn(An_j).
\end{equation}
In particular every denominator in the proof is nonzero, including those
at zero bases. Put $\delta_r=m_r-X_{i(r)}+(A-n_j)Y_{i(r)}$.
Direct expansion gives
\begin{equation}\label{eq:squareidentity}
 \sum_{r\in I_j}\delta_r^2
       =E_j-T_j\sum_{r\in I_j}Y_{i(r)}^2.
\end{equation}
Dividing \eqref{eq:R} by $\prod T_j$ shows that
$\sum_jn_jE_j/T_j$ is an integer, so
$\sum_j\sum_{r\in I_j}n_j\delta_r^2/T_j$ is an integer.

Let $J$ be the set of nonzero base groups, let $s=|J|\le5$, and define
\[
 C_0=\prod_{j\in J}n_j^b,\qquad
 d_r=C_0m_r-\frac{C_0}{n_j^b}n_j^{b+k_{i(r)}}\quad(j\in J).
\]
The empty product is one. In particular $C_0\ne0$; including a zero base
in this clearing product would invalidate the proof. The exponents and
quotients in this display define integers, and are proof notation only.
Lemma~\ref{lem:residues} gives $C_0\delta_r\equiv d_r\pmod{T_j}$.
Thus the rational number
\[
 V=\sum_{j\in J}\sum_{r\in I_j}\frac{n_jd_r^2}{T_j}
\]
is an integer. This follows by multiplying the previous integral rational
sum by $C_0^2$ and changing each numerator by a multiple of its own
denominator. It does not require different denominators to be coprime.

Since all bounded ports have absolute value at most $b$,
\[
 |d_r|\le b^{sb+1}+b^{(s+1)b}\le2b^{(s+1)b}.
\]
Hence
\[
 \sum_{j\in J}\sum_{r\in I_j}|n_j|d_r^2
 \le68b^{12b+1}<W^{6b+2}\le W^W.
\]
Here $68<W$, $b^2<W$, and $6b+2\le b^2+1\le W$ because $b\ge16$.
Together with \eqref{eq:Tbounds}, this gives $|V|<1$. Every summand has
the common sign $\sgn(A)$ or is zero, because $n_j/T_j$ has that sign.
The integer $V$ is therefore zero, and each $d_r$ is zero. Divide by the
nonzero $C_0$ in $\mathbb Q$ to obtain the claimed powers. If $|n_j|\ge2$
and $k=-u<0$, the resulting equation $m_rn_j^u=1$ is impossible over
$\Z$.

Conversely, exact powers at a nonzero base give $T_j\mid\delta_r$ and
then $T_j\mid E_j$ by \eqref{eq:squareidentity}. The integer
$z=\sum_{j\in J}n_jE_j/T_j$ supplies \eqref{eq:R}. Zero bases contribute
zero. This proves both directions.
\end{proof}

\section{Authenticating indices and powers together}
\label{sec:joint}

The larger gain in proof structure is that canonicality need not be
established in a separate divisibility certificate. Let
\begin{equation}\label{eq:Omega}
 L=A-1,\qquad
 \Omega=\sum_{i=1}^{14}\bigl((X_i-1)^2+(Y_i-k_i)^2\bigr).
\end{equation}
The sum is over distinct pairs, not over the seventeen requests; a pair
used at several bases occurs only once.

\begin{theorem}[Joint index and power authentication]\label{thm:joint}
Assume the bounds \eqref{eq:bounds}--\eqref{eq:largeA}, the norm and
$Z_i=AY_i$ equations, and the raw classification bound \eqref{eq:rawclass}.
No correctness of the requested indices is assumed. There is an integer
$z$ satisfying the single polynomial condition
\begin{equation}\label{eq:G}
 \boxed{\mathcal G=
 \Omega\prod_jT_j+
 L\sum_j n_jE_j\prod_{\ell\ne j}T_\ell
       -zL\prod_jT_j=0}
\end{equation}
if and only if all fourteen pairs are canonical at their requested signed
indices and every request at a nonzero base is its exact power value.
\end{theorem}

\begin{proof}
The denominator $L$ has sign $\sgn(A)$ and $|L|\ge W^W$. For the raw
representation \eqref{eq:rawclass}, Lemma~\ref{lem:residues} gives
\[
 \Omega\equiv o:=\sum_i\bigl((\varepsilon_i-1)^2+
                     (\varepsilon_i\nu_i-k_i)^2\bigr)\pmod L.
\]
The integer $o$ is nonnegative. For a nonzero base, the raw power residue is
\[
 X_i+(n_j-A)Y_i\equiv\varepsilon_i n_j^{\nu_i}\pmod{T_j}.
\]
The exponent here is the classified $\nu_i$, whose bound is already
known, rather than the as-yet unauthenticated $k_i$.

As before, let $C_0$ be the product of $n_j^b$ over the nonzero base
groups, and put
\[
 d_r=C_0m_r-\varepsilon_{i(r)}\frac{C_0}{n_j^b}n_j^{b+\nu_{i(r)}}.
\]
Dividing \eqref{eq:G} by its nonzero common denominator, using
\eqref{eq:squareidentity}, multiplying by $C_0^2$, and taking the small
residue representatives shows that
\begin{equation}\label{eq:jointsmall}
 V=\frac{C_0^2o}{L}+
       \sum_{j:n_j\ne0}\sum_{r\in I_j}\frac{n_jd_r^2}{T_j}
       \quad\hbox{is an integer}.
\end{equation}
Every summand has the common sign $\sgn(A)$ or is zero. Sharing factors
between $C_0$, $L$, and the $T_j$ causes no difficulty: only integrality
of a rational sum and congruent numerators have been used.

Bound fourteen by fifteen in the elementary estimate
$o\le60(1+b^2)\le120b^2$. Then
\begin{align*}
 C_0^2o+\sum |n_j|d_r^2
 &\le120b^{10b+2}+68b^{12b+1}\\
 &\le188b^{12b+2}<W^{6b+2}\le W^W.
\end{align*}
The strict inequality uses $188<W$ and $b^2<W$. Every active denominator
in \eqref{eq:jointsmall} has magnitude at least $W^W$, so $|V|<1$.
Consequently $V=0$, $o=0$, and every $d_r=0$.
The vanishing of $o$ gives $\varepsilon_i=1$ and $\nu_i=k_i$ for every
pair. The other vanishing errors give the exact powers at precisely those
authenticated indices. This proves soundness without presupposing
canonicality.

For completeness, canonical pairs make $L\mid\Omega$, and exact
nonzero-base powers make $T_j\mid E_j$. Choose
\[
 z=\frac{\Omega}{L}+\sum_{j:n_j\ne0}\frac{n_jE_j}{T_j}.
\]
This is an integer and satisfies \eqref{eq:G}.
\end{proof}

The rational expressions above are exclusively part of the proof.
Equation~\eqref{eq:G} is an ordinary polynomial with finite sums and
products. In particular it contains no variable exponent, division,
inequality predicate, or quantified coefficient parameter.

\subsection{Partial fusion and partitions}

One can partition the five base groups and give each part one quotient.
For a singleton part, use $E_j-z_jT_j=0$; the one-denominator argument
permits a zero base since then $T_j=-1$. For a larger part, use the
oriented numerator in \eqref{eq:R} restricted to that part. Each quotient
has a precise cost in variables, while its denominator product controls
degree.

The orientation term can be fused with just one part. If that part is
the growth base $t$, which is positive before any power interpretation,
use the particularly economical residual
\begin{equation}\label{eq:tfusion}
 \Omega T_t+L(E_t-z_tT_t)=0.
\end{equation}
Here $1/T_t$ and $1/L$ have the same sign. The proof of
Theorem~\ref{thm:joint}, with the clearing factor $t^b$, authenticates
all fourteen indices and the growth power. The remaining parts can then
use Theorem~\ref{thm:power}. Thus the family includes a 210-variable
degree-at-most-eight seed with five singleton power parts, and a
208-variable degree-at-most-twelve seed with parts
$\{t\},\{B,BH\},\{q,q+1\}$ and fusion at $t$.

For consecutive bases $q,q+1$, a two-base part may use unit weights:
the denominators attached to nonzero bases have a common sign. At $q=0$ or $q=-1$,
the zero-base contribution has denominator $-1$ and can be removed as
an integer before the small-error argument. This is another valid
partition choice; it does not assume $q>0$ prematurely.

\section{The literal consistency polynomial and its bootstrap}
\label{sec:literal}

Let $Q_0,\ldots,Q_{152}$ be the exact parent quadratic residuals.
All retained ones are printed in Appendix~\ref{app:rows}; their coefficients
are also in \code{source/equation216.json}. For the main fused construction,
delete the coordinates
\[
 141,142,143,144,145,146,201,212,213,214,215
\]
and the rows
\[
 97,98,99,100,101,102,142,149,150,151,152.
\]
The ten power coordinates occur only in the ten specified power rows,
and coordinate 201 occurs only in row 142. The remaining 205 coordinates
are independently quantified integers; add one integer $z$. Define
\begin{equation}\label{eq:P206}
 \boxed{P_{206}=\mathcal G^2+
  \sum_{r\in\{0,\ldots,96\}\cup\{103,\ldots,141\}\cup\{143,\ldots,148\}}
  Q_r^2.}
\end{equation}
This is a complete finite definition, using the literal port table in
Appendix~\ref{app:ports} and \eqref{eq:T}--\eqref{eq:G}.
The JSON circuit renumbers the retained coordinates increasingly and
places $z$ last; original indices in the paper are names, not additional
unknowns.

The exact degree is 26. The residual $\mathcal R$ has degree at most 12,
and the monomial
\[
 x_0^2x_4^2x_{94}^4x_{130}x_{184}x_{186}^2
\]
has coefficient $-2$. Its unique degree-12 contribution involving this
output comes from the request with output $x_{130}$ in the $q+1$ group.
Multiplying by the $x_{134}$ term of $L$ gives the degree-13 monomial
\[
 x_0^2x_4^2x_{94}^4x_{130}x_{134}x_{184}x_{186}^2
\]
with coefficient $-2$ in $\mathcal G$. The other summands in
\eqref{eq:G} have degree at most 12. The square of this nonzero leading
homogeneous polynomial is nonzero, while all other squares in
\eqref{eq:P206} have degree at most four. The independent checker also
computes this coefficient exactly in a truncated polynomial ring.

The intermediate $P_{207}$ retains coordinate 201 and row 142, replaces
only the ten power coordinates and rows, and uses $\mathcal R^2$ in
place of $\mathcal G^2$. It has exact degree 24.

\subsection{Why the zero-base clause is sufficient}

The five actual bases are
\[
 t=x_0,\quad B=x_4,\quad BH=x_{186},\quad q=x_{94}-1,
 \quad q+1=x_{94}.
\]
Neither $B$ nor $q$ may be assumed nonzero at the start. The joint
theorem first gives canonical indices and exactness at whichever bases
are nonzero. This is enough for the following acyclic deductions.

\begin{enumerate}
\item The horizon $t=16+4d^2\ge16$ is nonzero. Its request gives
 $F=t^t$, so $F$ is positive and even, with $F\ge2^{64}$.
\item The retained row $Bh_B=C_{\rm crt}+F+1$ has an odd right side,
 since $C_{\rm crt}=2c$. Thus $B$ is odd and nonzero. The actual
 $H=2+4(C_{\rm crt}+a_1+\cdots+a_8)$ is 2 modulo 4 for any signed
 representatives, so $|BH|\ge2$.
\item The request at $BH$ gives $x\ge0$ and $q=(BH)^x$. The beta
 spacing is $y=F(2v+1)$, hence $|y|\ge F$. Set $L_0=1+(t+1)y$;
 then $|L_0|>1$. The terminal endpoint is $x-1=L_0u_t$.
 If $x=0$ it would make $L_0$ divide $-1$, a contradiction.
 Thus $x\ge1$ and $|q|\ge2$.
\item The $q$ request with index $y+1$ forces $y+1\ge0$.
 A negative $y$ would satisfy $y\le-F<-1$, so $y\ge F>0$.
\item From $x=F+1+(y+1)u_0$, a negative $u_0$ would imply
 $x\le F-y\le0$. Hence $u_0\ge0$ and $x\ge F+1$.
 The terminal endpoint now gives $u_t\ge1$ and
 $x\ge2+(t+1)y$. Thus $|q|\ge2^x$ and $q+1$ is a nonzero nonunit.
\item Every request is therefore exact at a nonzero base. The unchanged
 signed binomial-digit argument authenticates
 $B=\binom{F+t}{t}>0$ and recovers the nonnegative CRT representatives.
 It then gives $H>0$ and $q>0$. The inherited local range and
 positive-weight proof-flow decoding produces a finite C222 refutation.
\end{enumerate}

All elementary rows in this schedule were independently recovered as
combinations of retained rows. The last step imports the stated parent
binomial/CRT and proof-calculus theorems; it is not inferred from numerical
testing or from an unspecified appeal to DPRM.

\subsection{Restoring the parent and composing the logical calibration}

After this bootstrap all bases are nonzero. For each old singleton,
restore its modulus as $T_j$ and its quotient as $-\delta_r/T_j$.
For each old multi-request pool, restore
\[
 L_j=E_j/T_j,\qquad K_j=2H_j+(2A-n_j)L_j.
\]
For the orientation row, restore $Q=\Omega/(A-1)$.
These divisions yield integers by the theorems above. The orientation
row in the centered source is not literally $\Omega-(A-1)Q$; its
recorded combination with retained norm and frame rows is. Since those
other rows vanish, the restored quotient makes the actual row vanish.
Thus a new zero extends to an actual zero of the original $P_{216}$.

Conversely a parent zero supplies canonical pairs and exact powers, so
the displayed formulas supply the new quotient $z$. This proves an
existential equivalence over $\Z$; the restorations need not be polynomial
maps on arbitrary tuples. Composing it with the inherited declarations
\begin{quote}\small
\code{KEndpoint.c222_iff_zf_refutation}\\
\code{ConstructibleChoice.zf_refutation_iff_zfc_refutation}
\end{quote}
and the parent arithmetic decoding proves \eqref{eq:target} for $P_{206}$,
$P_{207}$, and the partition seeds.

\section{Removing defined arithmetic coordinates}
\label{sec:elimination}

The semantic compression above can be followed by an unconditional
algebraic transformation. This second method is inherited from source
152; the new work applies it to the new literal seeds and records the
resulting equations and certificates.

\begin{lemma}[Integral monic substitution]\label{lem:monic}
Suppose one residual is $\epsilon x_v+f(\widehat{\boldsymbol x})$,
where $\epsilon\in\{1,-1\}$ and $f$ does not contain $x_v$.
Delete that residual and substitute $x_v=-\epsilon f$ into every other
residual. Then the zero sets of the two systems are in bijection over
$\Z$, by coordinate deletion and the displayed polynomial restoration.
An acyclic sequence of such deletions has the same property.
\end{lemma}
\begin{proof}
At a zero, the selected equation uniquely forces $x_v=-\epsilon f$.
Substitution therefore preserves every other zero equation. Conversely,
a zero of the substituted system has a unique integer extension by this
formula. Restore an acyclic sequence in reverse order. No divisibility,
positivity, or existence theorem is needed.
\end{proof}

The optimizer selects a finite acyclic set of such source pivots and
computes a conservative degree bound for every substituted residual.
It does not claim to find the best possible pivot set. The delivered
certificate fixes the actual pivots, source identity, retained variables,
and substitution circuit, so its validity does not depend on trusting
the search heuristic.

Each final polynomial is the sum of the squares of its surviving
substituted residuals. A circuit gives an exact finite polynomial, even
when expanding every monomial would be inefficient. The only nodes are
literal integer constants, variable leaves, addition, subtraction, and
multiplication. Repeatedly substituting those definitions yields an
ordinary member of $\Z[\boldsymbol x]$.

An exact degree certificate consists of an upper bound propagated through
the circuit and a nonzero evaluation of the corresponding top homogeneous
part modulo a specified integer greater than one. This is a finite exact witness
that the homogeneous polynomial is not zero; it is not a probabilistic
identity test. Likewise, a nonzero partial-derivative value modulo an integer greater than one
certifies genuine dependence on a coordinate. The independent verifier
checks the certificates as well as the source pivots and substitutions.

\section{Measured equations and their interpretation}
\label{sec:results}

\begin{center}\small
\begin{tabular}{rrrrrl}
\toprule
Variables & Degree & Residuals & Live operations & Height bound & Seed\\\midrule
194 & 6 & 131 & 2,682 & $2^{292}$ & $P_{211}$\\
147 & 8 & 84 & 2,559 & $2^{325}$ & $P_{211}$\\
135 & 12 & 72 & 2,498 & $2^{331}$ & $P_{211}$\\
121 & 16 & 58 & 2,453 & $2^{324}$ & $P_{211}$\\
115 & 24 & 52 & 2,438 & $2^{306}$ & $P_{209}$\\
111 & 32 & 48 & 2,429 & $2^{364}$ & $P_{209}$\\
105 & 64 & 42 & 2,495 & $2^{546}$ & $P_{207}$\\
97 & 256 & 34 & 2,398 & $2^{1670}$ & $P_{207}$\\
90 & 32,768 & 27 & 2,306 & $2^{114689}$ & $P_{206}$\\
\bottomrule\end{tabular}\end{center}
The height column is a checked upper bound on the maximum absolute expanded coefficient. 
The largest literal integer in each circuit is 384. Live operations count binary 
addition, subtraction, and multiplication in the emitted shared circuit, including squares.

The 90-variable equation has 27 residuals and exact degree 32,768. 
It is the smallest-arity equation produced in this package. The 121-variable 
degree-16 equation is a useful alternative when degree matters. Their exact files are 
\code{higher_degree/equations/integer90_degree32768.json} and 
\code{higher_degree/equations/integer121_degree16.json}.


All final variables range independently over unrestricted integers.
The displayed degrees are degrees of the resulting polynomial, rather
than of an expression with extra implicit witnesses. The literal
constants in the circuits remain small; expanded coefficients can grow
under substitution. Reported coefficient-size bounds, where given, are
proved upper bounds on the expanded coefficient height, not exact maxima.
The archive does not label a circuit's largest literal constant as the
height of its expanded polynomial.

The useful new comparisons are with the pinned direct-calculus branch.
Source 152, for example, supplies $(153,8)$, $(144,12)$, $(125,16)$,
$(117,32)$, $(111,64)$, $(102,256)$, and $(94,65540)$ in
(integer variables, degree). The present table identifies actual
improvements at several of these budgets. It does not claim to dominate
the repository's small-variable, large-coefficient constructions.

\section{Verification, reproducibility, and the proof boundary}
\label{sec:verification}

The archive supplies the full LaTeX source, compiled PDF, literal
polynomial circuits, retained residual lists, source ports, pivot
certificates, builders, evaluators, independent verifiers, and receipts
from actual successful runs. Python's standard library is sufficient
for the finite arithmetic checks. The final verification commands are
documented in the top-level README.

The validation is divided by what each layer can establish. The parent
replay authenticates the incoming finite construction. A separately
implemented port audit reconstructs the 36-square frame, all seventeen
request ports, all fourteen Pell pairs, the anchor comparisons, and
the exact bootstrap identities. It proves that their required row
combinations avoid the deleted rows. A separate circuit checker binds
the literal pooling expression and retained rows to those ports and
checks the degree certificate. A further verifier reconstructs every
monic substitution and final circuit from its certificate. Arithmetic
fixtures exercise negative parameters, signed indices, zero and unit
bases, repeated moduli, and deliberately incorrect signs and outputs.

These tasks are complementary. Finite fixtures do not prove the
unbounded theorem; source hashes do not prove the meaning of a
polynomial; and a numerical search that finds no zeros would not prove
consistency. The unbounded argument is the ordinary proof in this
article, composed with the precisely identified inherited arithmetic
and logical theorems. No new end-to-end Lean/Rocq build or encoded
object-language ZFC derivation was produced. The separate reviews were
AI-assisted internal checks, not external human peer review.

The construction is additive material for ProveIt intake. The public
repository was read at the pinned commit; this submission does not
require any modification to its existing files.

\section{A rejected structural shortcut: using the radix as the CRT modulus}
\label{sec:obstruction}

We also examined whether the distinct request $N=B^x$ could be removed by
using $q=(BH)^x$ as the CRT modulus. The unchanged range tests would then
require $q\mid\binom{C_{\rm crt}}{t}$. This seemingly economical
replacement cannot preserve the existing construction.

After the same sign and binomial recovery, all relevant representatives
are nonnegative,
\[
 B=\binom{F+t}{t}>F+t,\quad
 H>C_{\rm crt},\quad H>\max(1,C_{\rm crt}),\quad x\ge t.
\]
For $C_{\rm crt}\ge t$, one has
\[
 0<\binom{C_{\rm crt}}t
 \le C_{\rm crt}^t<H^x<(BH)^x=q.
\]
Thus the proposed divisibility would force $C_{\rm crt}<t$.
But the retained row requires $B\mid C_{\rm crt}+F+1$, while
\[
 0<C_{\rm crt}+F+1\le F+t<B,
\]
a contradiction. Therefore that specific substitution makes the
unchanged range-and-index system inconsistent. A successful removal of
the separate CRT modulus would have to redesign these range conditions,
not merely identify the two existing coordinates.

This obstruction helps explain the source architecture: the radix must
dominate coefficients for digit recovery, whereas the CRT modulus must
permit large global representatives with independently prescribed
small local residues. Those two requirements cannot be merged by the
proposed literal substitution.

\section{What remains open}

The joint certificate removes the separation between index authentication
and exact-power recovery, but retains the squared anchor that bounds
the classified indices. A further reduction would need either a cheaper
way to establish that raw bound, or a new simultaneous certificate that
also authenticates the anchor without a circular size argument. The
present proof does not provide such a replacement.

A second concrete optimization is to select the partition and monic
eliminations jointly. A larger pool uses fewer quotient variables but
raises the degree of some frequently used ports, which can make later
substitution more expensive. The measured family demonstrates that
pooling every block is not automatically the best choice at a fixed
degree. Exhaustive global optimization of this design space was not
performed.

Finally, formalizing the rational-integrality lemma and its exact source
binding would provide a comparatively small new proof-assistant target.
Its essential ingredients are congruence modulo the integers, a common
sign for all rational error terms, and a strict upper bound below one.
The large inherited logical calibration need not be redesigned to
formalize this local improvement.

\appendix
\section{Complete affine port table}\label{app:ports}

The source coordinates are $x_0,\ldots,x_{215}$. Set $A=x_{134}+1$ and
use $D_i=X_i-AY_i$ as the stored Pell difference. This difference port
is named \code{D} in the electronic data. Each abscissa is recovered as
$X_i=D_i+Z_i$.
The bases indexed $0,\ldots,4$ are
$(x_0,x_4,x_{186},x_{94}-1,x_{94})$.
The explicit derived-label index is
\begin{align*}
 h={}&-4-3x_8-7x_9+4x_{19}+2x_{22}+2x_{26}+2x_{28}+2x_{30}+2x_{32}\\
 &+4x_{34}+2x_{35}+2x_{38}+2x_{41}+4x_{60}+2x_{61}.
\end{align*}
\begin{center}\small
\begin{longtable}{rclllll}
\toprule Request & Base group & Output $m$ & Index $k$ & Difference $D_i$ & $Y_i$ & $Z_i$\\\midrule\endhead
0 & 0 & $x_{3}$ & $x_{0}$ & $x_{138}$ & $x_{139}$ & $x_{140}$\\
1 & 1 & $x_{17}$ & $x_{1}$ & $x_{174}$ & $x_{175}$ & $1 +x_{176}$\\
2 & 2 & $-1 +x_{94}$ & $x_{1}$ & $x_{174}$ & $x_{175}$ & $1 +x_{176}$\\
3 & 4 & $x_{95}$ & $x_{0} +x_{3}$ & $x_{147}$ & $x_{148}$ & $x_{149}$\\
4 & 3 & $x_{96}$ & $x_{0}$ & $x_{138}$ & $x_{139}$ & $x_{140}$\\
5 & 4 & $x_{100}$ & $-1 +x_{8}$ & $x_{150}$ & $-1 +x_{151}$ & $x_{152}$\\
7 & 4 & $x_{103}$ & $x_{9}$ & $x_{153}$ & $x_{154}$ & $x_{155}$\\
8 & 4 & $x_{106}$ & $x_{10}$ & $x_{156}$ & $x_{157}$ & $x_{158}$\\
9 & 4 & $x_{109}$ & $x_{11}$ & $x_{159}$ & $x_{160}$ & $x_{161}$\\
10 & 4 & $x_{112}$ & $x_{12}$ & $x_{162}$ & $x_{163}$ & $x_{164}$\\
11 & 3 & $x_{113}$ & $x_{2}$ & $x_{165}$ & $x_{166}$ & $x_{167}$\\
12 & 4 & $x_{116}$ & $x_{13}$ & $x_{168}$ & $x_{169}$ & $x_{170}$\\
13 & 4 & $x_{119}$ & $x_{14}$ & $x_{171}$ & $x_{172}$ & $x_{173}$\\
14 & 3 & $x_{120}$ & $x_{1}$ & $x_{174}$ & $x_{175}$ & $1 +x_{176}$\\
15 & 4 & $x_{124}$ & $x_{15}$ & $x_{177}$ & $x_{178}$ & $x_{179}$\\
16 & 4 & $x_{127}$ & $2x_{16}$ & $x_{180}$ & $x_{181}$ & $x_{182}$\\
17 & 4 & $x_{130}$ & $h$ & $x_{183}$ & $x_{184}$ & $x_{185}$\\
\bottomrule\end{longtable}\end{center}

The fourteen distinct pairs in $\Omega$ are represented once each by
requests 0, 3, 5, 7, 8, 9, 10, 11, 12, 13, 14, 15, 16, and 17.
All symbols in this table are explicit affine polynomials.

\section{Every retained quadratic residual}\label{app:rows}

The following list prints all 143 retained rows of the intermediate
$P_{207}$. For the joint construction $P_{206}$ omit the one row labeled
$Q_{142}$ as specified in \eqref{eq:P206}. No listed row contains any of
the ten deleted power coordinates. Row $Q_{142}$ is the only listed
row containing coordinate $x_{201}$. The authoritative machine-readable
integer coefficient list and an exact parser check accompany this display.

{\footnotesize
% BEGIN_RESIDUAL 0
\begin{align*}
 Q_{0}={}&x_{0} -x_{0}x_{2} -x_{2} +x_{7}
\end{align*}
% END_RESIDUAL 0
% BEGIN_RESIDUAL 1
\begin{align*}
 Q_{1}={}&-x_{2}x_{5} -x_{3} +x_{6}x_{7}
\end{align*}
% END_RESIDUAL 1
% BEGIN_RESIDUAL 2
\begin{align*}
 Q_{2}={}&-1 +x_{1} -x_{6}x_{7}
\end{align*}
% END_RESIDUAL 2
% BEGIN_RESIDUAL 3
\begin{align*}
 Q_{3}={}&1 +x_{3} -x_{4}x_{18} +2x_{16}
\end{align*}
% END_RESIDUAL 3
% BEGIN_RESIDUAL 4
\begin{align*}
 Q_{4}={}&-2 -x_{8} -x_{8}^{2} -2x_{8}x_{9} -2x_{9} -x_{9}^{2} +x_{19}
\end{align*}
% END_RESIDUAL 4
% BEGIN_RESIDUAL 5
\begin{align*}
 Q_{5}={}&-12 +4x_{8} -2x_{8}^{2} -8x_{8}x_{20} -8x_{20}^{2} +x_{21} +x_{33}
\end{align*}
% END_RESIDUAL 5
% BEGIN_RESIDUAL 6
\begin{align*}
 Q_{6}={}&44 -90x_{8} +12x_{8}^{2} +32x_{8}x_{19} -24x_{8}x_{21} -384x_{19} -44x_{19}^{2} -192x_{19}x_{20} -128x_{19}x_{21} -229x_{20}\\
 &{}-64x_{20}^{2} -64x_{20}x_{21} -48x_{21} -4x_{21}^{2} +x_{22} +2x_{33}
\end{align*}
% END_RESIDUAL 6
% BEGIN_RESIDUAL 7
\begin{align*}
 Q_{7}={}&-2 -x_{9}^{2} -2x_{9}x_{10} -x_{10}^{2} +x_{23}
\end{align*}
% END_RESIDUAL 7
% BEGIN_RESIDUAL 8
\begin{align*}
 Q_{8}={}&-1 +2x_{8} -x_{8}^{2} -2x_{8}x_{10} +2x_{10} -x_{10}^{2} +x_{24}
\end{align*}
% END_RESIDUAL 8
% BEGIN_RESIDUAL 9
\begin{align*}
 Q_{9}={}&3 -15x_{8} -9x_{8}^{2} -12x_{8}x_{9} -36x_{8}x_{10} -24x_{8}x_{19} -12x_{8}x_{23} -24x_{8}x_{24} +12x_{9} -28x_{10} -34x_{10}^{2}\\
 &{}-32x_{10}x_{19} -32x_{10}x_{24} +6x_{19} -8x_{19}x_{24} -x_{20} +12x_{23} -22x_{24} -4x_{24}^{2} +x_{25} -4x_{25}^{2} +x_{26} -3x_{43}
\end{align*}
% END_RESIDUAL 9
% BEGIN_RESIDUAL 10
\begin{align*}
 Q_{10}={}&-15 +14x_{8} -2x_{8}^{2} -8x_{8}x_{9} -8x_{8}x_{23} +4x_{9}^{2} +8x_{9}x_{23} +4x_{23}^{2} -4x_{25}^{2} +x_{26} -3x_{43}
\end{align*}
% END_RESIDUAL 10
% BEGIN_RESIDUAL 11
\begin{align*}
 Q_{11}={}&-6 +16x_{10} -16x_{10}^{2} -16x_{10}x_{23} -16x_{10}x_{24} +8x_{23} -4x_{23}^{2} -8x_{23}x_{24} +8x_{24} -4x_{24}^{2} +x_{27}
\end{align*}
% END_RESIDUAL 11
% BEGIN_RESIDUAL 12
\begin{align*}
 Q_{12}={}&-6 +23x_{10} +9x_{10}^{2} -16x_{10}x_{19} +8x_{10}x_{24} +4x_{10}x_{27} +10x_{19} -16x_{19}x_{24} -8x_{19}x_{27} -x_{20} +24x_{24}\\
 &{}+12x_{27} +x_{28} -x_{42}
\end{align*}
% END_RESIDUAL 12
% BEGIN_RESIDUAL 13
\begin{align*}
 Q_{13}={}&6 -19x_{8} -35x_{8}^{2} -80x_{8}x_{10} +4x_{8}x_{23} -24x_{8}x_{24} +33x_{9} +35x_{9}^{2} +80x_{9}x_{10} +24x_{9}x_{23} -4x_{9}x_{24}\\
 &{}+16x_{10} +32x_{10}x_{23} -32x_{10}x_{24} +12x_{23} +4x_{23}^{2} -8x_{24} -4x_{24}^{2} +x_{29} +4x_{29}^{2} -x_{30}
\end{align*}
% END_RESIDUAL 13
% BEGIN_RESIDUAL 14
\begin{align*}
 Q_{14}={}&90 -10x_{8} -2x_{8}^{2} -24x_{8}x_{9} -32x_{8}x_{10} -8x_{8}x_{23} -72x_{9} -18x_{9}^{2} -84x_{9}x_{10} -48x_{9}x_{23} -96x_{10}\\
 &{}-74x_{10}^{2} -64x_{10}x_{23} -78x_{23} -8x_{23}^{2} -4x_{29}^{2} +x_{30}
\end{align*}
% END_RESIDUAL 14
% BEGIN_RESIDUAL 15
\begin{align*}
 Q_{15}={}&24 -12x_{8} +20x_{8}^{2} +8x_{8}x_{9} +16x_{8}x_{10} +16x_{8}x_{19} +16x_{8}x_{23} -24x_{9} -36x_{9}^{2} +16x_{9}x_{10} +32x_{9}x_{19}\\
 &{}+32x_{9}x_{23} +24x_{10} -12x_{10}^{2} -8x_{10}x_{19} -8x_{10}x_{23} -12x_{19} -4x_{19}^{2} -8x_{19}x_{23} +32x_{23} -4x_{23}^{2} +x_{31}
\end{align*}
% END_RESIDUAL 15
% BEGIN_RESIDUAL 16
\begin{align*}
 Q_{16}={}&25 -56x_{8} -48x_{8}^{2} -128x_{8}x_{10} -64x_{8}x_{23} -32x_{8}x_{24} -32x_{8}x_{31} -90x_{10} -52x_{10}^{2} -48x_{10}x_{23}\\
 &{}-8x_{10}x_{24} -40x_{10}x_{31} -24x_{23} -8x_{23}^{2} -16x_{23}x_{31} -38x_{24} +4x_{24}^{2} -8x_{24}x_{31} -2x_{27} -4x_{31} -4x_{31}^{2}\\
 &{}+x_{32}
\end{align*}
% END_RESIDUAL 16
% BEGIN_RESIDUAL 17
\begin{align*}
 Q_{17}={}&1 -4x_{19}^{2} +x_{20}
\end{align*}
% END_RESIDUAL 17
% BEGIN_RESIDUAL 18
\begin{align*}
 Q_{18}={}&-6 +2x_{8} -x_{8}^{2} +8x_{8}x_{19} -4x_{8}x_{20} -8x_{19} +16x_{19}x_{20} -4x_{20}^{2} +x_{33}
\end{align*}
% END_RESIDUAL 18
% BEGIN_RESIDUAL 19
\begin{align*}
 Q_{19}={}&44 -90x_{8} +12x_{8}^{2} -32x_{8}x_{19} -24x_{8}x_{33} +384x_{19} -44x_{19}^{2} +192x_{19}x_{20} +128x_{19}x_{33} -229x_{20}\\
 &{}-64x_{20}^{2} -64x_{20}x_{33} +2x_{21} -48x_{33} -4x_{33}^{2} +x_{34}
\end{align*}
% END_RESIDUAL 19
% BEGIN_RESIDUAL 20
\begin{align*}
 Q_{20}={}&-x_{9}^{2} -4x_{9}x_{34} -4x_{34}^{2} +x_{35}
\end{align*}
% END_RESIDUAL 20
% BEGIN_RESIDUAL 21
\begin{align*}
 Q_{21}={}&34 +x_{8} +5x_{8}^{2} -22x_{8}x_{9} +32x_{8}x_{19} +2x_{9} -43x_{9}^{2} +64x_{9}x_{19} -3x_{19} -16x_{19}^{2} +x_{36}
\end{align*}
% END_RESIDUAL 21
% BEGIN_RESIDUAL 22
\begin{align*}
 Q_{22}={}&49 -34x_{9}^{2} -44x_{9}x_{10} -24x_{9}x_{23} -32x_{9}x_{36} +6x_{10} -10x_{10}^{2} -24x_{10}x_{23} -24x_{10}x_{36} -32x_{23}\\
 &{}-8x_{23}x_{36} +4x_{36} -4x_{36}^{2} +x_{37} -x_{43}
\end{align*}
% END_RESIDUAL 22
% BEGIN_RESIDUAL 23
\begin{align*}
 Q_{23}={}&8x_{8} -64x_{8}^{2} -96x_{8}x_{10} -32x_{8}x_{24} -32x_{8}x_{37} +6x_{10} -36x_{10}^{2} -24x_{10}x_{24} -24x_{10}x_{37} +2x_{24} -4x_{24}^{2}\\
 &{}-8x_{24}x_{37} +4x_{37} -4x_{37}^{2} +x_{38}
\end{align*}
% END_RESIDUAL 23
% BEGIN_RESIDUAL 24
\begin{align*}
 Q_{24}={}&-x_{10}^{2} -2x_{10}x_{11} -x_{11}^{2} +x_{39}
\end{align*}
% END_RESIDUAL 24
% BEGIN_RESIDUAL 25
\begin{align*}
 Q_{25}={}&-11 +3x_{8}^{2} +8x_{8}x_{10} -8x_{8}x_{24} -4x_{8}x_{39} +4x_{9}^{2} +8x_{9}x_{23} -8x_{9}x_{24} -8x_{10} +2x_{10}^{2} -8x_{10}x_{39} +4x_{23}^{2}\\
 &{}-8x_{23}x_{24} +18x_{24} -4x_{24}^{2} +4x_{25}^{2} -x_{26} +4x_{39} -4x_{39}^{2} +x_{40}
\end{align*}
% END_RESIDUAL 25
% BEGIN_RESIDUAL 26
\begin{align*}
 Q_{26}={}&4x_{8} -18x_{8}^{2} -40x_{8}x_{10} -32x_{8}x_{24} -40x_{8}x_{39} -16x_{8}x_{40} +4x_{10} -24x_{10}^{2} -32x_{10}x_{24} -48x_{10}x_{39}\\
 &{}-16x_{10}x_{40} -16x_{24}^{2} -32x_{24}x_{39} -16x_{24}x_{40} +4x_{39} -24x_{39}^{2} -16x_{39}x_{40} -4x_{40}^{2} +x_{41}
\end{align*}
% END_RESIDUAL 26
% BEGIN_RESIDUAL 27
\begin{align*}
 Q_{27}={}&x_{10} -25x_{10}^{2} -40x_{10}x_{24} -20x_{10}x_{27} -16x_{24}^{2} -16x_{24}x_{27} -4x_{27}^{2} +x_{42}
\end{align*}
% END_RESIDUAL 27
% BEGIN_RESIDUAL 28
\begin{align*}
 Q_{28}={}&-1 +8x_{9} -4x_{9}^{2} -8x_{9}x_{23} +8x_{23} -4x_{23}^{2} +x_{43}
\end{align*}
% END_RESIDUAL 28
% BEGIN_RESIDUAL 29
\begin{align*}
 Q_{29}={}&-1 -16x_{9}^{2} -8x_{9}x_{10} -32x_{9}x_{23} -16x_{9}x_{43} -x_{10}^{2} -8x_{10}x_{23} -4x_{10}x_{43} -16x_{23}^{2} -16x_{23}x_{43} -4x_{43}^{2}\\
 &{}+x_{44}
\end{align*}
% END_RESIDUAL 29
% BEGIN_RESIDUAL 30
\begin{align*}
 Q_{30}={}&3 +x_{9} -2x_{9}x_{10} -4x_{9}x_{44} -3x_{10}^{2} -8x_{10}x_{44} -x_{23} -4x_{44}^{2} +x_{45}
\end{align*}
% END_RESIDUAL 30
% BEGIN_RESIDUAL 31
\begin{align*}
 Q_{31}={}&-1 -4x_{45} -4x_{45}^{2} +x_{46}
\end{align*}
% END_RESIDUAL 31
% BEGIN_RESIDUAL 32
\begin{align*}
 Q_{32}={}&-1 -x_{10}^{2} -8x_{10}x_{45} -4x_{10}x_{46} -16x_{45}^{2} -16x_{45}x_{46} -4x_{46}^{2} +x_{47}
\end{align*}
% END_RESIDUAL 32
% BEGIN_RESIDUAL 33
\begin{align*}
 Q_{33}={}&4 +x_{9} -2x_{9}x_{10} -4x_{9}x_{47} -3x_{10}^{2} -8x_{10}x_{47} -x_{23} -4x_{47}^{2} +x_{48}
\end{align*}
% END_RESIDUAL 33
% BEGIN_RESIDUAL 34
\begin{align*}
 Q_{34}={}&-4x_{42}^{2} -8x_{42}x_{48} -4x_{48}^{2} +x_{49}
\end{align*}
% END_RESIDUAL 34
% BEGIN_RESIDUAL 35
\begin{align*}
 Q_{35}={}&-9 +7x_{8} -x_{8}^{2} -8x_{8}x_{42} -4x_{8}x_{49} +24x_{42} -16x_{42}^{2} -16x_{42}x_{49} +12x_{49} -4x_{49}^{2} +x_{50}
\end{align*}
% END_RESIDUAL 35
% BEGIN_RESIDUAL 36
\begin{align*}
 Q_{36}={}&-1 -4x_{50} -4x_{50}^{2} +x_{51}
\end{align*}
% END_RESIDUAL 36
% BEGIN_RESIDUAL 37
\begin{align*}
 Q_{37}={}&-x_{10}^{2} -8x_{10}x_{50} -4x_{10}x_{51} -16x_{50}^{2} -16x_{50}x_{51} -4x_{51}^{2} +x_{52}
\end{align*}
% END_RESIDUAL 37
% BEGIN_RESIDUAL 38
\begin{align*}
 Q_{38}={}&-x_{8} -x_{8}^{2} -4x_{8}x_{10} -4x_{8}x_{52} -4x_{10} -4x_{10}^{2} -8x_{10}x_{52} -4x_{52} -4x_{52}^{2} +x_{53}
\end{align*}
% END_RESIDUAL 38
% BEGIN_RESIDUAL 39
\begin{align*}
 Q_{39}={}&-10 -12x_{53} -4x_{53}^{2} +x_{54}
\end{align*}
% END_RESIDUAL 39
% BEGIN_RESIDUAL 40
\begin{align*}
 Q_{40}={}&2 -4x_{9} -16x_{9}x_{53} -8x_{9}x_{54} -4x_{23} -16x_{23}x_{53} -8x_{23}x_{54} -x_{43} +8x_{53} -16x_{53}^{2} -16x_{53}x_{54} +4x_{54}\\
 &{}-4x_{54}^{2} +x_{55}
\end{align*}
% END_RESIDUAL 40
% BEGIN_RESIDUAL 41
\begin{align*}
 Q_{41}={}&-10 +x_{9} -25x_{9}^{2} -40x_{9}x_{23} -20x_{9}x_{55} -16x_{23}^{2} -16x_{23}x_{55} -12x_{48} -4x_{48}^{2} -4x_{55}^{2} +x_{56}
\end{align*}
% END_RESIDUAL 41
% BEGIN_RESIDUAL 42
\begin{align*}
 Q_{42}={}&21 +24x_{48} -8x_{48}^{2} -16x_{48}x_{56} -2x_{56} -4x_{56}^{2} +x_{57}
\end{align*}
% END_RESIDUAL 42
% BEGIN_RESIDUAL 43
\begin{align*}
 Q_{43}={}&4 -4x_{8} -16x_{8}x_{19} +16x_{8}x_{59} -8x_{9} -32x_{9}x_{19} +32x_{9}x_{59} +x_{10} -x_{10}^{2} -4x_{10}x_{57} -6x_{19} -8x_{19}x_{59}\\
 &{}+3x_{20} -x_{36} -4x_{57}^{2} +x_{58} +12x_{59} -4x_{59}^{2}
\end{align*}
% END_RESIDUAL 43
% BEGIN_RESIDUAL 44
\begin{align*}
 Q_{44}={}&8 +12x_{8} +8x_{8}x_{10} -16x_{8}x_{19} +24x_{9} +16x_{9}x_{10} -32x_{9}x_{19} +3x_{10} -x_{10}^{2} -4x_{10}x_{19} -14x_{19}\\
 &{}+12x_{19}^{2} -x_{36} +x_{59}
\end{align*}
% END_RESIDUAL 44
% BEGIN_RESIDUAL 45
\begin{align*}
 Q_{45}={}&-28 -56x_{8} -48x_{8}x_{10} +32x_{8}x_{58} +32x_{8}x_{59} -112x_{9} -96x_{9}x_{10} +64x_{9}x_{58} +64x_{9}x_{59} -18x_{10}\\
 &{}+6x_{10}^{2} +24x_{10}x_{19} +28x_{19} -16x_{19}x_{58} -16x_{19}x_{59} +16x_{58} -4x_{58}^{2} +18x_{59} -8x_{59}^{2} +x_{60}
\end{align*}
% END_RESIDUAL 45
% BEGIN_RESIDUAL 46
\begin{align*}
 Q_{46}={}&-2 -4x_{60} -4x_{60}^{2} +x_{61}
\end{align*}
% END_RESIDUAL 46
% BEGIN_RESIDUAL 47
\begin{align*}
 Q_{47}={}&-x_{3}x_{13} +x_{62}
\end{align*}
% END_RESIDUAL 47
% BEGIN_RESIDUAL 48
\begin{align*}
 Q_{48}={}&x_{12} -x_{13} -2x_{22}x_{62} -x_{63}
\end{align*}
% END_RESIDUAL 48
% BEGIN_RESIDUAL 49
\begin{align*}
 Q_{49}={}&x_{12} -x_{13} -2x_{26}x_{62} -x_{64}
\end{align*}
% END_RESIDUAL 49
% BEGIN_RESIDUAL 50
\begin{align*}
 Q_{50}={}&-2x_{28}x_{62} +2x_{41}x_{62} -2x_{62} -x_{65} +x_{70}
\end{align*}
% END_RESIDUAL 50
% BEGIN_RESIDUAL 51
\begin{align*}
 Q_{51}={}&x_{12} -x_{13} -2x_{30}x_{62} -x_{66}
\end{align*}
% END_RESIDUAL 51
% BEGIN_RESIDUAL 52
\begin{align*}
 Q_{52}={}&x_{12} -x_{13} -2x_{32}x_{62} +6x_{62} -x_{67}
\end{align*}
% END_RESIDUAL 52
% BEGIN_RESIDUAL 53
\begin{align*}
 Q_{53}={}&x_{12} -x_{13} -4x_{34}x_{62} -2x_{35}x_{62} -x_{68}
\end{align*}
% END_RESIDUAL 53
% BEGIN_RESIDUAL 54
\begin{align*}
 Q_{54}={}&x_{12} -x_{13} -2x_{38}x_{62} -x_{69}
\end{align*}
% END_RESIDUAL 54
% BEGIN_RESIDUAL 55
\begin{align*}
 Q_{55}={}&x_{12} -x_{13} -2x_{41}x_{62} -4x_{62} -x_{70}
\end{align*}
% END_RESIDUAL 55
% BEGIN_RESIDUAL 56
\begin{align*}
 Q_{56}={}&-x_{8}x_{9} +x_{8}x_{10} +x_{9} +x_{9}x_{10} -x_{10} -x_{10}^{2} +x_{71}
\end{align*}
% END_RESIDUAL 56
% BEGIN_RESIDUAL 57
\begin{align*}
 Q_{57}={}&1 +2x_{11} +x_{11}^{2} -x_{71}^{2} +x_{72}
\end{align*}
% END_RESIDUAL 57
% BEGIN_RESIDUAL 58
\begin{align*}
 Q_{58}={}&-x_{12}^{2} +2x_{12}x_{13} +2x_{12}x_{74} -x_{13}^{2} -2x_{13}x_{74} -x_{72}^{2} +x_{73} -x_{74}^{2}
\end{align*}
% END_RESIDUAL 58
% BEGIN_RESIDUAL 59
\begin{align*}
 Q_{59}={}&-4x_{60}x_{62} -2x_{61}x_{62} +2x_{62} +x_{74}
\end{align*}
% END_RESIDUAL 59
% BEGIN_RESIDUAL 60
\begin{align*}
 Q_{60}={}&-x_{9}x_{62} -x_{13} +x_{75}x_{76} +x_{77}
\end{align*}
% END_RESIDUAL 60
% BEGIN_RESIDUAL 61
\begin{align*}
 Q_{61}={}&-1 +4x_{3} +4x_{3}x_{8} +8x_{3}x_{9} -2x_{3}x_{19} +x_{75}
\end{align*}
% END_RESIDUAL 61
% BEGIN_RESIDUAL 62
\begin{align*}
 Q_{62}={}&-x_{8}x_{211} -x_{12} +x_{76} +x_{211}
\end{align*}
% END_RESIDUAL 62
% BEGIN_RESIDUAL 63
\begin{align*}
 Q_{63}={}&x_{9}x_{211} +x_{12} -x_{13} -2x_{19}x_{62} +2x_{62} +x_{78}
\end{align*}
% END_RESIDUAL 63
% BEGIN_RESIDUAL 64
\begin{align*}
 Q_{64}={}&-x_{12}x_{63} +x_{13}x_{63} +x_{79}
\end{align*}
% END_RESIDUAL 64
% BEGIN_RESIDUAL 65
\begin{align*}
 Q_{65}={}&-x_{64}x_{79} +x_{80}
\end{align*}
% END_RESIDUAL 65
% BEGIN_RESIDUAL 66
\begin{align*}
 Q_{66}={}&-x_{65}x_{80} +x_{81}
\end{align*}
% END_RESIDUAL 66
% BEGIN_RESIDUAL 67
\begin{align*}
 Q_{67}={}&-x_{66}x_{81} +x_{82}
\end{align*}
% END_RESIDUAL 67
% BEGIN_RESIDUAL 68
\begin{align*}
 Q_{68}={}&-x_{67}x_{82} +x_{83}
\end{align*}
% END_RESIDUAL 68
% BEGIN_RESIDUAL 69
\begin{align*}
 Q_{69}={}&-x_{68}x_{83} +x_{84}
\end{align*}
% END_RESIDUAL 69
% BEGIN_RESIDUAL 70
\begin{align*}
 Q_{70}={}&-x_{69}x_{84} +x_{85}
\end{align*}
% END_RESIDUAL 70
% BEGIN_RESIDUAL 71
\begin{align*}
 Q_{71}={}&-x_{70}x_{85} +x_{86}
\end{align*}
% END_RESIDUAL 71
% BEGIN_RESIDUAL 72
\begin{align*}
 Q_{72}={}&-x_{77}x_{86} +x_{87}
\end{align*}
% END_RESIDUAL 72
% BEGIN_RESIDUAL 73
\begin{align*}
 Q_{73}={}&-x_{78}x_{87} +x_{88}
\end{align*}
% END_RESIDUAL 73
% BEGIN_RESIDUAL 74
\begin{align*}
 Q_{74}={}&-x_{88}^{2} +x_{89}
\end{align*}
% END_RESIDUAL 74
% BEGIN_RESIDUAL 75
\begin{align*}
 Q_{75}={}&-x_{2} -2x_{2}x_{16} +2x_{16} +x_{90}
\end{align*}
% END_RESIDUAL 75
% BEGIN_RESIDUAL 76
\begin{align*}
 Q_{76}={}&-x_{14}x_{90} +x_{91}
\end{align*}
% END_RESIDUAL 76
% BEGIN_RESIDUAL 77
\begin{align*}
 Q_{77}={}&-2x_{1}^{2} +2x_{1}x_{12} +2x_{1}x_{13} +2x_{1}x_{91} +2x_{1}x_{92} -x_{12}^{2} -2x_{12}x_{91} -x_{13}^{2} -2x_{13}x_{92} +x_{17}x_{93} -x_{73}x_{89}\\
 &{}-x_{91}^{2} -x_{92}^{2}
\end{align*}
% END_RESIDUAL 77
% BEGIN_RESIDUAL 78
\begin{align*}
 Q_{78}={}&-x_{2}x_{15} +x_{15} -x_{15}x_{90} +x_{92}
\end{align*}
% END_RESIDUAL 78
% BEGIN_RESIDUAL 79
\begin{align*}
 Q_{79}={}&-x_{94}x_{97} +x_{97} +x_{99}
\end{align*}
% END_RESIDUAL 79
% BEGIN_RESIDUAL 80
\begin{align*}
 Q_{80}={}&-x_{4}x_{120} +x_{95}x_{202} -x_{98} -x_{99}x_{120}
\end{align*}
% END_RESIDUAL 80
% BEGIN_RESIDUAL 81
\begin{align*}
 Q_{81}={}&x_{100}x_{202} -x_{101}x_{123} -x_{102}
\end{align*}
% END_RESIDUAL 81
% BEGIN_RESIDUAL 82
\begin{align*}
 Q_{82}={}&x_{103}x_{202} -x_{104}x_{123} -x_{105}
\end{align*}
% END_RESIDUAL 82
% BEGIN_RESIDUAL 83
\begin{align*}
 Q_{83}={}&x_{106}x_{202} -x_{107}x_{123} -x_{108}
\end{align*}
% END_RESIDUAL 83
% BEGIN_RESIDUAL 84
\begin{align*}
 Q_{84}={}&x_{109}x_{202} -x_{110}x_{123} -x_{111}
\end{align*}
% END_RESIDUAL 84
% BEGIN_RESIDUAL 85
\begin{align*}
 Q_{85}={}&x_{112}x_{203} -x_{114}x_{123} -x_{115}
\end{align*}
% END_RESIDUAL 85
% BEGIN_RESIDUAL 86
\begin{align*}
 Q_{86}={}&x_{116}x_{203} -x_{117}x_{123} -x_{118}
\end{align*}
% END_RESIDUAL 86
% BEGIN_RESIDUAL 87
\begin{align*}
 Q_{87}={}&-x_{17}x_{120} +x_{123}
\end{align*}
% END_RESIDUAL 87
% BEGIN_RESIDUAL 88
\begin{align*}
 Q_{88}={}&x_{119} -x_{121}x_{123} -x_{122}
\end{align*}
% END_RESIDUAL 88
% BEGIN_RESIDUAL 89
\begin{align*}
 Q_{89}={}&-x_{123}x_{125} +x_{124} -x_{126}
\end{align*}
% END_RESIDUAL 89
% BEGIN_RESIDUAL 90
\begin{align*}
 Q_{90}={}&-x_{123}x_{128} +x_{127}x_{202} -x_{129}
\end{align*}
% END_RESIDUAL 90
% BEGIN_RESIDUAL 91
\begin{align*}
 Q_{91}={}&-x_{123}x_{131} +x_{130}x_{202} -x_{132}
\end{align*}
% END_RESIDUAL 91
% BEGIN_RESIDUAL 92
\begin{align*}
 Q_{92}={}&x_{133} -2x_{133}x_{135} +x_{136}
\end{align*}
% END_RESIDUAL 92
% BEGIN_RESIDUAL 93
\begin{align*}
 Q_{93}={}&-x_{133}x_{136} -x_{136} +x_{137}
\end{align*}
% END_RESIDUAL 93
% BEGIN_RESIDUAL 94
\begin{align*}
 Q_{94}={}&2x_{134} +x_{134}^{2} +x_{136}^{2} -x_{137}^{2}
\end{align*}
% END_RESIDUAL 94
% BEGIN_RESIDUAL 95
\begin{align*}
 Q_{95}={}&-x_{134}x_{139} -x_{139} +x_{140}
\end{align*}
% END_RESIDUAL 95
% BEGIN_RESIDUAL 96
\begin{align*}
 Q_{96}={}&-1 +x_{138}^{2} +2x_{138}x_{140} +x_{139}^{2}
\end{align*}
% END_RESIDUAL 96
% BEGIN_RESIDUAL 103
\begin{align*}
 Q_{103}={}&-x_{134}x_{148} -x_{148} +x_{149}
\end{align*}
% END_RESIDUAL 103
% BEGIN_RESIDUAL 104
\begin{align*}
 Q_{104}={}&-1 +x_{147}^{2} +2x_{147}x_{149} +x_{148}^{2}
\end{align*}
% END_RESIDUAL 104
% BEGIN_RESIDUAL 105
\begin{align*}
 Q_{105}={}&1 +x_{134} -x_{134}x_{151} -x_{151} +x_{152}
\end{align*}
% END_RESIDUAL 105
% BEGIN_RESIDUAL 106
\begin{align*}
 Q_{106}={}&x_{150}^{2} +2x_{150}x_{152} -2x_{151} +x_{151}^{2}
\end{align*}
% END_RESIDUAL 106
% BEGIN_RESIDUAL 107
\begin{align*}
 Q_{107}={}&-x_{134}x_{154} -x_{154} +x_{155}
\end{align*}
% END_RESIDUAL 107
% BEGIN_RESIDUAL 108
\begin{align*}
 Q_{108}={}&-1 +x_{153}^{2} +2x_{153}x_{155} +x_{154}^{2}
\end{align*}
% END_RESIDUAL 108
% BEGIN_RESIDUAL 109
\begin{align*}
 Q_{109}={}&-x_{134}x_{157} -x_{157} +x_{158}
\end{align*}
% END_RESIDUAL 109
% BEGIN_RESIDUAL 110
\begin{align*}
 Q_{110}={}&-1 +x_{156}^{2} +2x_{156}x_{158} +x_{157}^{2}
\end{align*}
% END_RESIDUAL 110
% BEGIN_RESIDUAL 111
\begin{align*}
 Q_{111}={}&-x_{134}x_{160} -x_{160} +x_{161}
\end{align*}
% END_RESIDUAL 111
% BEGIN_RESIDUAL 112
\begin{align*}
 Q_{112}={}&-1 +x_{159}^{2} +2x_{159}x_{161} +x_{160}^{2}
\end{align*}
% END_RESIDUAL 112
% BEGIN_RESIDUAL 113
\begin{align*}
 Q_{113}={}&-x_{134}x_{163} -x_{163} +x_{164}
\end{align*}
% END_RESIDUAL 113
% BEGIN_RESIDUAL 114
\begin{align*}
 Q_{114}={}&-1 +x_{162}^{2} +2x_{162}x_{164} +x_{163}^{2}
\end{align*}
% END_RESIDUAL 114
% BEGIN_RESIDUAL 115
\begin{align*}
 Q_{115}={}&-x_{134}x_{166} -x_{166} +x_{167}
\end{align*}
% END_RESIDUAL 115
% BEGIN_RESIDUAL 116
\begin{align*}
 Q_{116}={}&-1 +x_{165}^{2} +2x_{165}x_{167} +x_{166}^{2}
\end{align*}
% END_RESIDUAL 116
% BEGIN_RESIDUAL 117
\begin{align*}
 Q_{117}={}&-x_{134}x_{169} -x_{169} +x_{170}
\end{align*}
% END_RESIDUAL 117
% BEGIN_RESIDUAL 118
\begin{align*}
 Q_{118}={}&-1 +x_{168}^{2} +2x_{168}x_{170} +x_{169}^{2}
\end{align*}
% END_RESIDUAL 118
% BEGIN_RESIDUAL 119
\begin{align*}
 Q_{119}={}&-x_{134}x_{172} -x_{172} +x_{173}
\end{align*}
% END_RESIDUAL 119
% BEGIN_RESIDUAL 120
\begin{align*}
 Q_{120}={}&-1 +x_{171}^{2} +2x_{171}x_{173} +x_{172}^{2}
\end{align*}
% END_RESIDUAL 120
% BEGIN_RESIDUAL 121
\begin{align*}
 Q_{121}={}&1 -x_{134}x_{175} -x_{175} +x_{176}
\end{align*}
% END_RESIDUAL 121
% BEGIN_RESIDUAL 122
\begin{align*}
 Q_{122}={}&-1 +2x_{174} +x_{174}^{2} +2x_{174}x_{176} +x_{175}^{2}
\end{align*}
% END_RESIDUAL 122
% BEGIN_RESIDUAL 123
\begin{align*}
 Q_{123}={}&-x_{134}x_{178} -x_{178} +x_{179}
\end{align*}
% END_RESIDUAL 123
% BEGIN_RESIDUAL 124
\begin{align*}
 Q_{124}={}&-1 +x_{177}^{2} +2x_{177}x_{179} +x_{178}^{2}
\end{align*}
% END_RESIDUAL 124
% BEGIN_RESIDUAL 125
\begin{align*}
 Q_{125}={}&-x_{134}x_{181} -x_{181} +x_{182}
\end{align*}
% END_RESIDUAL 125
% BEGIN_RESIDUAL 126
\begin{align*}
 Q_{126}={}&-1 +x_{180}^{2} +2x_{180}x_{182} +x_{181}^{2}
\end{align*}
% END_RESIDUAL 126
% BEGIN_RESIDUAL 127
\begin{align*}
 Q_{127}={}&-x_{134}x_{184} -x_{184} +x_{185}
\end{align*}
% END_RESIDUAL 127
% BEGIN_RESIDUAL 128
\begin{align*}
 Q_{128}={}&-1 +x_{183}^{2} +2x_{183}x_{185} +x_{184}^{2}
\end{align*}
% END_RESIDUAL 128
% BEGIN_RESIDUAL 129
\begin{align*}
 Q_{129}={}&-1 -2x_{0} -4x_{0}^{2} -2x_{0}x_{3} -x_{1}^{2} -x_{2}^{2} -2x_{3}^{2} -x_{4}^{2} +2x_{8} -x_{8}^{2} +2x_{9}x_{10} -x_{11}^{2} -x_{12}^{2} -x_{13}^{2} -x_{14}^{2}\\
 &{}-x_{15}^{2} -4x_{16}^{2} -x_{17}^{2} -x_{23} +2x_{94} -2x_{94}^{2} -x_{95}^{2} -x_{96}^{2} -x_{100}^{2} -x_{103}^{2} -x_{106}^{2} -x_{109}^{2} -x_{112}^{2} -x_{113}^{2}\\
 &{}-x_{116}^{2} -x_{119}^{2} -x_{120}^{2} -x_{124}^{2} -x_{127}^{2} -x_{130}^{2} +x_{133} +x_{139}^{2} +x_{148}^{2} -2x_{151} +x_{151}^{2} +x_{154}^{2} +x_{157}^{2}\\
 &{}+x_{160}^{2} +x_{163}^{2} +x_{166}^{2} +x_{169}^{2} +x_{172}^{2} +x_{175}^{2} +x_{178}^{2} +x_{181}^{2} -x_{184}^{2} -x_{186}^{2} +x_{204}^{2} +x_{205}^{2} +x_{206}\\
 &{}+x_{206}^{2}
\end{align*}
% END_RESIDUAL 129
% BEGIN_RESIDUAL 130
\begin{align*}
 Q_{130}={}&2x_{4} -4x_{4}x_{8} -4x_{4}x_{9} -4x_{4}x_{10} -4x_{4}x_{11} -4x_{4}x_{12} -4x_{4}x_{13} -4x_{4}x_{14} -4x_{4}x_{15} -8x_{4}x_{16} +x_{186}
\end{align*}
% END_RESIDUAL 130
% BEGIN_RESIDUAL 131
\begin{align*}
 Q_{131}={}&-16 +x_{0} -4x_{187}^{2}
\end{align*}
% END_RESIDUAL 131
% BEGIN_RESIDUAL 132
\begin{align*}
 Q_{132}={}&-x_{134}x_{189} -x_{189} +x_{198}
\end{align*}
% END_RESIDUAL 132
% BEGIN_RESIDUAL 133
\begin{align*}
 Q_{133}={}&-1 +x_{188}^{2} +x_{189}^{2} -x_{198}^{2}
\end{align*}
% END_RESIDUAL 133
% BEGIN_RESIDUAL 134
\begin{align*}
 Q_{134}={}&-x_{191} -x_{191}x_{192} +x_{199}
\end{align*}
% END_RESIDUAL 134
% BEGIN_RESIDUAL 135
\begin{align*}
 Q_{135}={}&-1 +x_{190}^{2} +x_{191}^{2} -x_{199}^{2}
\end{align*}
% END_RESIDUAL 135
% BEGIN_RESIDUAL 136
\begin{align*}
 Q_{136}={}&-x_{184}^{2} +x_{200}
\end{align*}
% END_RESIDUAL 136
% BEGIN_RESIDUAL 137
\begin{align*}
 Q_{137}={}&x_{189} -2x_{194}x_{200} -x_{200}
\end{align*}
% END_RESIDUAL 137
% BEGIN_RESIDUAL 138
\begin{align*}
 Q_{138}={}&-4x_{189} -8x_{189}x_{193} +x_{192}
\end{align*}
% END_RESIDUAL 138
% BEGIN_RESIDUAL 139
\begin{align*}
 Q_{139}={}&-x_{134} -x_{188}x_{195} +x_{192}
\end{align*}
% END_RESIDUAL 139
% BEGIN_RESIDUAL 140
\begin{align*}
 Q_{140}={}&-x_{183} -x_{185} -x_{188}x_{196} +x_{190}
\end{align*}
% END_RESIDUAL 140
% BEGIN_RESIDUAL 141
\begin{align*}
 Q_{141}={}&4 +3x_{8} +7x_{9} -4x_{19} -2x_{22} -2x_{26} -2x_{28} -2x_{30} -2x_{32} -4x_{34} -2x_{35} -2x_{38} -2x_{41} -4x_{60} -2x_{61}\\
 &{}-4x_{184}x_{197} +x_{191}
\end{align*}
% END_RESIDUAL 141
% BEGIN_RESIDUAL 142
\begin{align*}
 Q_{142}={}&21 -2x_{0} -2x_{0}^{2} -2x_{0}x_{139} -2x_{0}x_{148} -2x_{1}x_{175} -2x_{2}x_{166} -x_{3}^{2} -2x_{3}x_{148} -x_{4}^{2} +2x_{8} -2x_{8}x_{151}\\
 &{}+6x_{8}x_{184} -2x_{9}x_{154} +14x_{9}x_{184} -2x_{10}x_{157} -2x_{11}x_{160} -2x_{12}x_{163} -2x_{13}x_{169} -2x_{14}x_{172} -2x_{15}x_{178}\\
 &{}-4x_{16}x_{181} -x_{17}^{2} -8x_{19}x_{184} -4x_{22}x_{184} -4x_{26}x_{184} -4x_{28}x_{184} -4x_{30}x_{184} -4x_{32}x_{184} -8x_{34}x_{184}\\
 &{}-4x_{35}x_{184} -4x_{38}x_{184} -4x_{41}x_{184} -8x_{60}x_{184} -4x_{61}x_{184} +2x_{94} -2x_{94}^{2} -x_{95}^{2} -x_{96}^{2} -x_{100}^{2} -x_{103}^{2}\\
 &{}-x_{106}^{2} -x_{109}^{2} -x_{112}^{2} -x_{113}^{2} -x_{116}^{2} -x_{119}^{2} -x_{120}^{2} -x_{124}^{2} -x_{127}^{2} -x_{130}^{2} +x_{133} -x_{134}x_{201}\\
 &{}-2x_{138} -2x_{140} +x_{140}^{2} -2x_{147} -2x_{149} +x_{149}^{2} -2x_{150} +2x_{151} -2x_{152} +x_{152}^{2} -2x_{153} -2x_{155} +x_{155}^{2}\\
 &{}-2x_{156} -2x_{158} +x_{158}^{2} -2x_{159} -2x_{161} +x_{161}^{2} -2x_{162} -2x_{164} +x_{164}^{2} -2x_{165} -2x_{167} +x_{167}^{2} -2x_{168}\\
 &{}-2x_{170} +x_{170}^{2} -2x_{171} -2x_{173} +x_{173}^{2} -2x_{174} +x_{176}^{2} -2x_{177} -2x_{179} +x_{179}^{2} -2x_{180} -2x_{182} +x_{182}^{2}\\
 &{}-2x_{183} +8x_{184} -2x_{185} +x_{185}^{2} -x_{186}^{2}
\end{align*}
% END_RESIDUAL 142
% BEGIN_RESIDUAL 143
\begin{align*}
 Q_{143}={}&-x_{96}x_{202} +x_{120}
\end{align*}
% END_RESIDUAL 143
% BEGIN_RESIDUAL 144
\begin{align*}
 Q_{144}={}&-x_{113}x_{203} +x_{120}
\end{align*}
% END_RESIDUAL 144
% BEGIN_RESIDUAL 145
\begin{align*}
 Q_{145}={}&-1 +x_{2} -x_{3} -2x_{3}x_{210}
\end{align*}
% END_RESIDUAL 145
% BEGIN_RESIDUAL 146
\begin{align*}
 Q_{146}={}&-10 -3x_{8} +8x_{8}^{2} +24x_{8}x_{9} -8x_{8}x_{19} +12x_{8}x_{22} +12x_{8}x_{26} +12x_{8}x_{28} +12x_{8}x_{30} +12x_{8}x_{32}\\
 &{}+24x_{8}x_{34} +12x_{8}x_{35} +12x_{8}x_{38} +12x_{8}x_{41} +24x_{8}x_{60} +12x_{8}x_{61} -14x_{9} +11x_{9}^{2} -8x_{9}x_{19} +28x_{9}x_{22}\\
 &{}+28x_{9}x_{26} +28x_{9}x_{28} +28x_{9}x_{30} +28x_{9}x_{32} +36x_{9}x_{34} +28x_{9}x_{35} +28x_{9}x_{38} +28x_{9}x_{41} +56x_{9}x_{60}\\
 &{}+28x_{9}x_{61} +13x_{19} -16x_{19}x_{22} -16x_{19}x_{26} -16x_{19}x_{28} -16x_{19}x_{30} -16x_{19}x_{32} -32x_{19}x_{34} -16x_{19}x_{35}\\
 &{}-16x_{19}x_{38} -16x_{19}x_{41} -32x_{19}x_{60} -16x_{19}x_{61} +16x_{22} -4x_{22}^{2} -8x_{22}x_{26} -8x_{22}x_{28} -8x_{22}x_{30}\\
 &{}-8x_{22}x_{32} -16x_{22}x_{34} -8x_{22}x_{35} -8x_{22}x_{38} -8x_{22}x_{41} -16x_{22}x_{60} -8x_{22}x_{61} +16x_{26} -4x_{26}^{2} -8x_{26}x_{28}\\
 &{}-8x_{26}x_{30} -8x_{26}x_{32} -16x_{26}x_{34} -8x_{26}x_{35} -8x_{26}x_{38} -8x_{26}x_{41} -16x_{26}x_{60} -8x_{26}x_{61} +16x_{28} -4x_{28}^{2}\\
 &{}-8x_{28}x_{30} -8x_{28}x_{32} -16x_{28}x_{34} -8x_{28}x_{35} -8x_{28}x_{38} -8x_{28}x_{41} -16x_{28}x_{60} -8x_{28}x_{61} +16x_{30} -4x_{30}^{2}\\
 &{}-8x_{30}x_{32} -16x_{30}x_{34} -8x_{30}x_{35} -8x_{30}x_{38} -8x_{30}x_{41} -16x_{30}x_{60} -8x_{30}x_{61} +16x_{32} -4x_{32}^{2} -16x_{32}x_{34}\\
 &{}-8x_{32}x_{35} -8x_{32}x_{38} -8x_{32}x_{41} -16x_{32}x_{60} -8x_{32}x_{61} +32x_{34} -36x_{34}^{2} -16x_{34}x_{35} -16x_{34}x_{38}\\
 &{}-16x_{34}x_{41} -32x_{34}x_{60} -16x_{34}x_{61} +21x_{35} -4x_{35}^{2} -8x_{35}x_{38} -8x_{35}x_{41} -16x_{35}x_{60} -8x_{35}x_{61} -x_{36}\\
 &{}+16x_{38} -4x_{38}^{2} -8x_{38}x_{41} -16x_{38}x_{60} -8x_{38}x_{61} +16x_{41} -4x_{41}^{2} -16x_{41}x_{60} -8x_{41}x_{61} +28x_{60} -20x_{60}^{2}\\
 &{}-16x_{60}x_{61} +17x_{61} -4x_{61}^{2} -x_{139}^{2} -x_{148}^{2} +2x_{151} -x_{151}^{2} -x_{154}^{2} -x_{157}^{2} -x_{160}^{2} -x_{163}^{2} -x_{166}^{2}\\
 &{}-x_{169}^{2} -x_{172}^{2} -x_{175}^{2} -x_{178}^{2} -x_{181}^{2} +x_{184}^{2} -x_{204}^{2} -x_{205}^{2} -x_{206} -x_{206}^{2}
\end{align*}
% END_RESIDUAL 146
% BEGIN_RESIDUAL 147
\begin{align*}
 Q_{147}={}&-1 -4x_{98}^{2} -4x_{102}^{2} -4x_{105}^{2} -4x_{108}^{2} -4x_{111}^{2} -4x_{115}^{2} -4x_{118}^{2} +x_{120}^{2} -4x_{122}^{2} -4x_{126}^{2} -4x_{129}^{2} -4x_{132}^{2} -x_{207}^{2}\\
 &{}-x_{208}^{2} -x_{209} -x_{209}^{2}
\end{align*}
% END_RESIDUAL 147
% BEGIN_RESIDUAL 148
\begin{align*}
 Q_{148}={}&-x_{3}x_{12} +x_{211}
\end{align*}
% END_RESIDUAL 148
}


\section{Finite elimination specifications}\label{app:pivots}

Each listed pair $(r,v)$ selects the source row $r$ as a monic definition
of source coordinate $v$. The exact sign, polynomial, dependency order,
and original-to-final coordinate map are in the accompanying JSON
certificate; the verifier reconstructs them from the literal source
rows. Each final expression is obtained by the substitutions of
Lemma~\ref{lem:monic}, followed by a sum of the surviving squares.
The source seed and all row indices are explicitly identified for each
profile.

\subsection*{194 variables, degree 6}
Seed: \code{source211.json}. Coordinates and row labels below are original P216 indices; the pool quotients have synthetic labels beginning at 216.
{\small\raggedright
$(49,64)$, $(70,85)$, $(62,76)$, $(58,73)$, $(50,65)$, $(72,87)$, $(52,67)$, $(139,192)$, $(0,7)$, $(142,133)$, $(75,90)$, $(51,66)$, $(79,99)$, $(53,68)$, $(48,63)$, $(136,200)$, $(87,123)$.\par}

\subsection*{147 variables, degree 8}
Seed: \code{source211.json}. Coordinates and row labels below are original P216 indices; the pool quotients have synthetic labels beginning at 216.
{\small\raggedright
$(89,126)$, $(65,80)$, $(66,81)$, $(61,75)$, $(140,190)$, $(82,105)$, $(52,67)$, $(85,115)$, $(58,73)$, $(54,69)$, $(62,76)$, $(91,132)$, $(68,83)$, $(81,102)$, $(93,137)$, $(56,71)$, $(63,78)$, $(132,198)$, $(70,85)$, $(88,122)$, $(136,200)$, $(41,56)$, $(141,191)$, $(74,89)$, $(111,161)$, $(121,176)$, $(83,108)$, $(48,63)$, $(127,185)$, $(78,92)$, $(80,98)$, $(107,155)$, $(76,91)$, $(125,182)$, $(0,7)$, $(84,111)$, $(86,118)$, $(113,164)$, $(72,87)$, $(90,129)$, $(29,44)$, $(105,152)$, $(22,37)$, $(35,50)$, $(139,192)$, $(115,167)$, $(103,149)$, $(119,173)$, $(32,47)$, $(109,158)$, $(95,140)$, $(39,54)$, $(117,170)$, $(5,33)$, $(123,179)$, $(43,58)$, $(25,40)$, $(15,31)$, $(12,28)$, $(10,26)$, $(20,35)$, $(6,22)$, $(46,61)$, $(14,30)$.\par}

\subsection*{135 variables, degree 12}
Seed: \code{source211.json}. Coordinates and row labels below are original P216 indices; the pool quotients have synthetic labels beginning at 216.
{\small\raggedright
$(72,87)$, $(83,108)$, $(89,126)$, $(68,83)$, $(49,64)$, $(64,79)$, $(79,99)$, $(71,86)$, $(74,89)$, $(54,69)$, $(50,65)$, $(51,66)$, $(62,76)$, $(56,71)$, $(63,78)$, $(93,137)$, $(140,190)$, $(134,199)$, $(84,111)$, $(61,75)$, $(81,102)$, $(139,192)$, $(69,84)$, $(80,98)$, $(85,115)$, $(90,129)$, $(130,186)$, $(132,198)$, $(66,81)$, $(119,173)$, $(58,73)$, $(0,7)$, $(88,122)$, $(123,179)$, $(136,200)$, $(75,90)$, $(29,44)$, $(82,105)$, $(48,63)$, $(42,57)$, $(121,176)$, $(53,68)$, $(115,167)$, $(91,132)$, $(86,118)$, $(111,161)$, $(105,152)$, $(76,91)$, $(113,164)$, $(34,49)$, $(37,52)$, $(117,170)$, $(109,158)$, $(55,70)$, $(40,55)$, $(107,155)$, $(59,74)$, $(78,92)$, $(22,37)$, $(95,140)$, $(32,47)$, $(92,136)$, $(125,182)$, $(5,21)$, $(103,149)$, $(127,185)$, $(25,40)$, $(11,27)$, $(44,59)$, $(2,1)$, $(15,31)$, $(18,33)$, $(87,123)$, $(14,30)$, $(9,26)$, $(1,3)$.\par}

\subsection*{121 variables, degree 16}
Seed: \code{source211.json}. Coordinates and row labels below are original P216 indices; the pool quotients have synthetic labels beginning at 216.
{\small\raggedright
$(52,67)$, $(72,87)$, $(66,81)$, $(49,64)$, $(61,75)$, $(55,70)$, $(139,192)$, $(136,200)$, $(74,89)$, $(48,63)$, $(53,68)$, $(79,99)$, $(54,69)$, $(51,66)$, $(63,78)$, $(58,73)$, $(76,91)$, $(78,92)$, $(62,76)$, $(50,65)$, $(70,85)$, $(64,79)$, $(68,83)$, $(60,77)$, $(22,37)$, $(25,40)$, $(15,31)$, $(81,102)$, $(90,129)$, $(82,105)$, $(83,108)$, $(88,122)$, $(85,115)$, $(89,126)$, $(84,111)$, $(91,132)$, $(86,118)$, $(80,98)$, $(87,123)$, $(5,33)$, $(132,198)$, $(0,7)$, $(36,51)$, $(32,47)$, $(42,57)$, $(93,137)$, $(140,190)$, $(141,191)$, $(40,55)$, $(34,49)$, $(56,71)$, $(29,44)$, $(44,59)$, $(92,136)$, $(59,74)$, $(75,90)$, $(134,199)$, $(43,58)$, $(11,27)$, $(105,152)$, $(109,158)$, $(127,185)$, $(125,182)$, $(117,170)$, $(123,179)$, $(107,155)$, $(103,149)$, $(119,173)$, $(111,161)$, $(113,164)$, $(115,167)$, $(95,140)$, $(121,176)$, $(130,186)$, $(2,1)$, $(144,120)$, $(10,26)$, $(13,30)$, $(20,35)$, $(6,22)$, $(46,61)$, $(24,39)$, $(21,36)$, $(37,52)$, $(28,43)$, $(31,46)$, $(39,54)$, $(3,3)$, $(131,0)$, $(137,189)$.\par}

\subsection*{115 variables, degree 24}
Seed: \code{source209.json}. Coordinates and row labels below are original P216 indices; the pool quotients have synthetic labels beginning at 216.
{\small\raggedright
$(52,67)$, $(64,79)$, $(51,66)$, $(65,80)$, $(66,81)$, $(132,198)$, $(67,82)$, $(50,65)$, $(89,126)$, $(70,85)$, $(63,78)$, $(58,73)$, $(93,137)$, $(69,84)$, $(140,190)$, $(62,76)$, $(73,88)$, $(42,57)$, $(60,77)$, $(83,108)$, $(84,111)$, $(0,7)$, $(53,68)$, $(91,132)$, $(57,72)$, $(130,186)$, $(49,64)$, $(81,102)$, $(55,70)$, $(86,118)$, $(71,86)$, $(56,71)$, $(134,199)$, $(40,55)$, $(32,47)$, $(95,140)$, $(88,119)$, $(29,44)$, $(80,98)$, $(125,182)$, $(107,155)$, $(121,176)$, $(75,90)$, $(82,105)$, $(31,46)$, $(103,149)$, $(48,63)$, $(90,129)$, $(36,51)$, $(54,69)$, $(79,99)$, $(85,115)$, $(61,75)$, $(119,173)$, $(109,158)$, $(139,192)$, $(92,136)$, $(34,49)$, $(136,200)$, $(22,37)$, $(115,167)$, $(59,74)$, $(127,185)$, $(39,54)$, $(43,58)$, $(123,179)$, $(111,161)$, $(141,191)$, $(148,211)$, $(113,164)$, $(76,91)$, $(78,92)$, $(144,120)$, $(15,31)$, $(105,152)$, $(137,189)$, $(117,170)$, $(5,21)$, $(18,33)$, $(25,40)$, $(11,27)$, $(12,42)$, $(37,52)$, $(44,59)$, $(21,36)$, $(2,1)$, $(24,39)$, $(145,2)$, $(28,43)$, $(131,0)$, $(10,26)$, $(87,123)$, $(9,20)$, $(46,61)$.\par}

\subsection*{111 variables, degree 32}
Seed: \code{source209.json}. Coordinates and row labels below are original P216 indices; the pool quotients have synthetic labels beginning at 216.
{\small\raggedright
$(62,76)$, $(65,80)$, $(71,86)$, $(52,67)$, $(66,81)$, $(84,111)$, $(58,73)$, $(60,77)$, $(54,69)$, $(90,129)$, $(49,64)$, $(64,79)$, $(81,102)$, $(91,132)$, $(85,115)$, $(72,87)$, $(63,78)$, $(132,198)$, $(50,65)$, $(79,99)$, $(56,71)$, $(70,85)$, $(88,122)$, $(140,190)$, $(51,66)$, $(67,82)$, $(80,98)$, $(53,68)$, $(93,137)$, $(89,126)$, $(134,199)$, $(61,75)$, $(103,149)$, $(42,57)$, $(37,52)$, $(29,44)$, $(74,89)$, $(31,46)$, $(86,118)$, $(83,108)$, $(82,105)$, $(57,72)$, $(76,91)$, $(0,7)$, $(109,158)$, $(68,83)$, $(92,136)$, $(137,189)$, $(136,200)$, $(130,186)$, $(30,45)$, $(75,90)$, $(41,56)$, $(36,51)$, $(138,192)$, $(125,182)$, $(107,155)$, $(111,161)$, $(48,63)$, $(78,92)$, $(115,167)$, $(55,70)$, $(121,176)$, $(12,42)$, $(117,170)$, $(40,55)$, $(148,211)$, $(95,140)$, $(143,120)$, $(123,179)$, $(141,191)$, $(127,185)$, $(34,49)$, $(44,36)$, $(25,40)$, $(105,152)$, $(119,173)$, $(59,74)$, $(113,164)$, $(22,37)$, $(5,21)$, $(145,2)$, $(33,48)$, $(11,27)$, $(15,31)$, $(131,0)$, $(87,123)$, $(3,3)$, $(2,1)$, $(20,35)$, $(18,33)$, $(14,30)$, $(10,26)$, $(17,20)$, $(24,39)$, $(45,60)$, $(46,61)$, $(139,134)$.\par}

\subsection*{105 variables, degree 64}
Seed: \code{source207.json}. Coordinates and row labels below are original P216 indices; the pool quotients have synthetic labels beginning at 216.
{\small\raggedright
$(74,89)$, $(89,126)$, $(56,71)$, $(62,76)$, $(93,137)$, $(80,98)$, $(90,129)$, $(81,102)$, $(69,84)$, $(55,70)$, $(86,118)$, $(83,108)$, $(65,80)$, $(51,66)$, $(64,79)$, $(54,69)$, $(79,99)$, $(140,190)$, $(58,73)$, $(134,199)$, $(60,77)$, $(68,83)$, $(66,81)$, $(63,78)$, $(132,198)$, $(85,115)$, $(61,75)$, $(111,161)$, $(88,122)$, $(67,82)$, $(52,67)$, $(48,63)$, $(109,158)$, $(127,185)$, $(50,65)$, $(71,86)$, $(138,192)$, $(70,85)$, $(57,72)$, $(31,46)$, $(42,57)$, $(91,132)$, $(73,88)$, $(53,68)$, $(49,64)$, $(84,111)$, $(82,105)$, $(37,52)$, $(136,200)$, $(129,133)$, $(39,54)$, $(105,152)$, $(113,164)$, $(76,91)$, $(130,186)$, $(75,90)$, $(78,92)$, $(123,179)$, $(5,21)$, $(30,45)$, $(59,74)$, $(141,191)$, $(34,49)$, $(119,173)$, $(43,58)$, $(117,170)$, $(41,56)$, $(121,176)$, $(103,149)$, $(0,7)$, $(92,136)$, $(38,53)$, $(115,167)$, $(15,31)$, $(125,182)$, $(21,36)$, $(95,140)$, $(25,40)$, $(36,51)$, $(107,155)$, $(32,47)$, $(29,44)$, $(18,33)$, $(148,211)$, $(11,27)$, $(22,37)$, $(1,3)$, $(137,189)$, $(87,123)$, $(44,59)$, $(143,120)$, $(2,1)$, $(12,42)$, $(9,20)$, $(139,134)$, $(46,61)$, $(24,39)$, $(20,35)$, $(10,26)$, $(14,30)$, $(16,32)$, $(8,24)$.\par}

\subsection*{97 variables, degree 256}
Seed: \code{source207.json}. Coordinates and row labels below are original P216 indices; the pool quotients have synthetic labels beginning at 216.
{\small\raggedright
$(0,7)$, $(1,3)$, $(2,1)$, $(4,19)$, $(5,21)$, $(6,22)$, $(7,23)$, $(8,24)$, $(9,20)$, $(10,26)$, $(11,27)$, $(12,28)$, $(13,30)$, $(15,31)$, $(16,32)$, $(18,33)$, $(19,34)$, $(20,35)$, $(21,36)$, $(22,37)$, $(23,38)$, $(24,39)$, $(25,40)$, $(26,41)$, $(27,42)$, $(28,43)$, $(29,44)$, $(30,45)$, $(31,46)$, $(32,47)$, $(34,49)$, $(35,50)$, $(36,51)$, $(38,53)$, $(39,54)$, $(40,55)$, $(41,56)$, $(42,57)$, $(43,58)$, $(44,59)$, $(46,61)$, $(47,62)$, $(48,12)$, $(49,64)$, $(50,65)$, $(51,66)$, $(52,67)$, $(53,68)$, $(54,69)$, $(55,70)$, $(56,71)$, $(57,72)$, $(58,73)$, $(59,74)$, $(60,77)$, $(61,75)$, $(62,76)$, $(63,78)$, $(64,79)$, $(65,80)$, $(66,81)$, $(67,82)$, $(69,84)$, $(70,85)$, $(71,86)$, $(72,87)$, $(74,89)$, $(75,90)$, $(76,91)$, $(78,92)$, $(79,99)$, $(80,98)$, $(81,102)$, $(82,105)$, $(83,108)$, $(84,111)$, $(85,115)$, $(86,118)$, $(87,123)$, $(88,119)$, $(89,124)$, $(90,129)$, $(91,132)$, $(92,136)$, $(93,137)$, $(95,140)$, $(103,149)$, $(105,152)$, $(107,155)$, $(109,158)$, $(111,161)$, $(113,164)$, $(115,167)$, $(117,170)$, $(119,173)$, $(121,176)$, $(123,179)$, $(125,182)$, $(127,185)$, $(130,186)$, $(131,0)$, $(132,198)$, $(134,199)$, $(136,200)$, $(137,189)$, $(138,192)$, $(139,134)$, $(140,183)$, $(141,191)$, $(143,120)$.\par}

\subsection*{90 variables, degree 32,768}
Seed: \code{source206_joint.json}. Coordinates and row labels below are original P216 indices; the pool quotients have synthetic labels beginning at 216.
{\small\raggedright
$(65,80)$, $(66,81)$, $(67,82)$, $(68,83)$, $(69,84)$, $(70,85)$, $(71,86)$, $(72,87)$, $(73,88)$, $(74,89)$, $(136,200)$, $(76,91)$, $(148,211)$, $(143,120)$, $(87,123)$, $(47,62)$, $(64,79)$, $(79,99)$, $(81,102)$, $(82,105)$, $(83,108)$, $(84,111)$, $(85,115)$, $(86,118)$, $(88,122)$, $(89,126)$, $(90,129)$, $(91,132)$, $(93,137)$, $(132,198)$, $(134,199)$, $(138,192)$, $(92,136)$, $(95,140)$, $(103,149)$, $(107,155)$, $(109,158)$, $(111,161)$, $(113,164)$, $(115,167)$, $(117,170)$, $(119,173)$, $(123,179)$, $(125,182)$, $(127,185)$, $(137,189)$, $(2,1)$, $(131,0)$, $(1,3)$, $(17,20)$, $(139,134)$, $(49,64)$, $(51,66)$, $(54,69)$, $(60,77)$, $(62,76)$, $(80,98)$, $(140,190)$, $(0,7)$, $(48,63)$, $(75,90)$, $(121,176)$, $(31,46)$, $(36,51)$, $(59,74)$, $(78,92)$, $(34,49)$, $(39,54)$, $(24,39)$, $(20,35)$, $(46,61)$, $(50,65)$, $(52,67)$, $(53,68)$, $(57,72)$, $(55,70)$, $(105,152)$, $(7,23)$, $(61,75)$, $(63,78)$, $(56,71)$, $(42,57)$, $(37,52)$, $(5,21)$, $(28,43)$, $(8,24)$, $(4,19)$, $(58,73)$, $(32,47)$, $(27,42)$, $(30,45)$, $(18,33)$, $(33,48)$, $(38,53)$, $(130,186)$, $(41,56)$, $(35,50)$, $(11,27)$, $(21,36)$, $(10,26)$, $(29,44)$, $(44,59)$, $(12,28)$, $(22,37)$, $(23,38)$, $(6,22)$, $(19,34)$, $(141,191)$, $(14,30)$, $(43,58)$, $(26,41)$, $(25,40)$, $(45,60)$, $(15,31)$, $(16,32)$, $(129,133)$.\par}



\begin{thebibliography}{9}
\bibitem{parent} ProveIt incoming research archive,
\emph{Direct quartic consistency equations in 218 and 216 integers},
25 October 2026. Immutable comparison:
\url{https://github.com/VladimirReshetnikov/ProveIt/blob/315154b670d7bdb1679554734e417410c9b8187b/docs/incoming/ConZFC_Pooled_Projection_2026-10-25.zip}.
The complete original archive is retained in \code{source/}.
\bibitem{pool168} ProveIt report 168,
\emph{Reciprocal pooling and integral projections for Con(ZFC)},
25 October 2026. Pinned path:
\url{https://github.com/VladimirReshetnikov/ProveIt/tree/315154b670d7bdb1679554734e417410c9b8187b/Computability/DiophantineConsistency/con-zfc-diophantine-forms/168-reciprocal-pooling}.
\bibitem{elim152} ProveIt report 152,
\emph{Lucas corrections and elimination of arithmetic witnesses},
25 October 2026. Pinned path:
\url{https://github.com/VladimirReshetnikov/ProveIt/tree/315154b670d7bdb1679554734e417410c9b8187b/Computability/DiophantineConsistency/con-zfc-diophantine-forms/152-lucas-corrections}.
\bibitem{anchor135} ProveIt report 135,
\emph{Recovering nonnegative arithmetic from signed Pell orientation},
section ``A signed nominal index in the squared anchor.''
\code{135-cofinal-anchors-round17/direct/ORIENTATION_FIRST_PROOF.md},
with its retained \code{PROOF_315.md} dependencies; pinned to the same
comparison commit. Relevant source proofs are retained in this submission.
\bibitem{bayer} J.~Bayer, M.~David, M.~Ha\ss ler, Y.~Matiyasevich,
and D.~Schleicher, \emph{Diophantine equations over $\Z$: universal bounds
and parallel formalization}, arXiv:2506.20909v3, 22 October 2026.
\url{https://arxiv.org/abs/2506.20909}.
\end{thebibliography}
\end{document}
